\PassOptionsToPackage{dvipsnames,svgnames,table}{xcolor}
\documentclass[10pt]{article}
\usepackage[a4paper,margin=22mm]{geometry}
\usepackage[no-math]{fontspec}\setmainfont{Latin Modern Roman}
\usepackage{amsmath,amssymb,amsthm,mathtools,graphicx,xcolor,hyperref,xurl,enumitem,ifthen}
\setlength{\emergencystretch}{4em}
\SetMathAlphabet{\mathbf}{normal}{OT1}{cmr}{bx}{n}
\SetMathAlphabet{\mathbf}{bold}{OT1}{cmr}{bx}{n}
\newtheorem{theorem}{Theorem}[section]
\newtheorem{lemma}[theorem]{Lemma}
\newtheorem{proposition}[theorem]{Proposition}
\newtheorem{corollary}[theorem]{Corollary}
\newtheorem{conjecture}[theorem]{Conjecture}
\theoremstyle{definition}
\newtheorem{definition}[theorem]{Definition}
\newtheorem{example}[theorem]{Example}
\newtheorem{remark}[theorem]{Remark}
\newtheorem{notation}[theorem]{Notation}
% Strict exact1-only installed renderer; original hardcoded triv alias below is unchanged.
\newfontfamily\CorpusTrivialOne[Path=/usr/share/fonts/truetype/noto/]{NotoSansMath-Regular.ttf}
\DeclareRobustCommand{\mathbbm}[1]{\ifthenelse{\equal{#1}{1}}{\mathord{\text{\upshape\CorpusTrivialOne\char"1D7D9}}}{\PackageError{alco262}{Uninventoried mathbbm argument}{Only exact literal1 is authorized}}}
\newcommand{\triv}{\mathbbm{1}}
\newcommand{\Ch}{\operatorname{Char}}
\newcommand{\Infl}{\operatorname{Infl}}
\newcommand{\Irr}{\operatorname{Irr}}
\newcommand{\sgn}{\operatorname{sgn}}
\newcommand{\Syl}{\operatorname{Syl}}
\newcommand{\down}{\big\downarrow}
\newcommand{\up}{\big\uparrow}
\newcommand{\cB}{\mathcal{B}}
\newcommand{\cL}{\mathcal{L}}
\newcommand{\cP}{\mathcal{P}}
\newcommand{\cX}{\mathcal{X}}
\newcommand{\ff}{\mathsf{f}}
\newcommand{\fg}{\mathsf{g}}
\newcommand{\sk}{\mathsf{k}}
\newcommand{\sm}{\mathsf{m}}
\newcommand{\sn}{\mathsf{N}}
\newcommand{\bC}{\mathbb{C}}
\newcommand{\N}{\mathbb{N}}
\newcommand{\Z}{\mathbb{Z}}

%\equalenv{example}{exam}
%\equalenv{definition}{defi}
%\equalenv{notation}{nota}
%\equalenv{remark}{rema}
%\equalenv{proposition}{prop}
%\equalenv{corollary}{coro}
%
\newtheorem*{thmA}{Theorem A}
\newtheorem*{thmB}{Theorem B}
\newtheorem*{thmC}{Theorem C}

\allowdisplaybreaks
\begin{document}
\section*{1108. Alternating skew-plethysm recursion with one-row inner partition and output length n minus k}
Stacey Law and Yuji Okitani. \emph{On plethysms and Sylow branching coefficients}, Algebraic Combinatorics 6(2) (2023), 321--357. \url{https://doi.org/10.5802/alco.262}. Original official publisher version; CC BY 4.0, \url{https://creativecommons.org/licenses/by/4.0/}.\par
\paragraph{Selected problem (editorial).} Prove the exact author Theorem5.1 (the Section5 restatement of TheoremA): all stated natural-number, partition, output-length and skew-coefficient conventions remain as written. Complete proof863--912, the authored m=1 case914, full new skew/Mackey/operator core918--1074 and terminal m-at-least2 proof1076--1082 are retained. TheoremB is supporting construction within this packet, not a separate counted unit. Unrelated Sylow outcomes are not selected.
\paragraph{Literal source context (editorial).} Native1020 prints the second factor with subscript m, beside the explicit m-minus1 expression1013 and authored final conclusion1022. The whole Mackey calculation980--1023 remains literal. Native1065 prints inflation to \(T_w\) before induction from \(U\), beside1063's exact subgroup definition and the full1065--1066/1070--1073 transitivity chain. Native869 displays the Littlewood--Richardson coefficient predicate without an explicit equality, beside867's explicit zero/one criterion and complete bijection/equality872. The final1076--1082 proof uses the conjugate-indexed formulation alongside the exact858--861 statement. Every original token and surrounding author argument is retained; no factor, inflation target, predicate or conjugation is replaced and no editor proof is supplied.
\paragraph{Finite typography (editorial).} Exactly73 inventoried literal trivial-character aliases expand the author's hardcoded \texttt{mathbbm\{1\}} through pinned installed Noto U+1D7D9 as a mathematical ordinary atom, with proportional normal/script styles and every original exponent/subscript untouched. The double-struck-one category is preserved; sans-serif stroke, width and placement differ from the original Type3serif glyph. Every other argument is rejected. No general bbm package is emulated or acquired. Unused tableau/list/table imports remain in archived native evidence; no such controls occur in the selected ranges. Only exact outside-excerpt reference indexes are rendered as their printed source numbers; standalone styles/page flow differ. All statements, proofs and general mathematical expressions are editable native TeX.
\clearpage
\setcounter{section}{1}
% BEGIN NATIVE AUTHOR BLOCK 126--339
\section{Preliminaries}\label{sec:prelims}
Throughout, we use $\N$ to denote the set of natural numbers, and $\N_0$ for the set of non-negative integers. 
As stated in the introduction, for a finite group $G$ we use $\Irr(G)$ to denote a complete set of the ordinary irreducible characters of $G$. Further, we use $\Ch(G)$ to denote the set of all ordinary characters of $G$, and $\triv_G$ to mean the trivial character of $G$ (omitting the subscript when the meaning is clear from context).
For a subgroup $H\le G$ and $\phi\in\Irr(H)$, we let $\Irr(G\mid \phi)$ denote the set of $\chi\in\Irr(G)$ such that the restriction of $\chi$ to $H$ contains $\phi$ as a constituent. 

For $g\in G$ and $H\le G$, let $H^g\coloneqq gHg^{-1}$. Given $\phi\in\Ch(H)$, the character $\phi_g\in\Ch(H^g)$ is defined by $\phi_g(x)=\phi(g^{-1}xg)$. Mackey's Theorem allows us to describe restrictions and inductions between  subgroups of a finite group (see \cite[Chapter 5]{I}, for example).
\begin{theorem}[Mackey]\label{thm:Mackey}
	Let $G$ be a finite group and $H,K\le G$. Let $\phi\in\Ch(H)$. Then
	\[ \phi\up_H^G\down_K = \sum_{g\in K\setminus G/H} \phi_g \up_{H^g\cap K}^K, \]
	where the sum runs over a set of $(K,H)$-double coset representatives.
\end{theorem}



\subsection{Representation theory of symmetric groups}\label{sec:prelims-sn}
Next, we recall some key facts concerning the representation theory of symmetric groups, and refer the reader to \cite{J,JK} for further detail.
It is well known that $\Irr(S_n)$ is naturally in bijection with the set $\cP(n)$ of all partitions of $n$. By convention, $\cP(0)=\{\emptyset\}$ where $\emptyset$ denotes the empty partition. 
The irreducible character of $S_n$ corresponding to the partition $\lambda\vdash n$ will be denoted by $\chi^\lambda$. In particular, $\chi^{(n)}=\triv_{S_n}$, the trivial character of $S_n$, and $\chi^{(1^n)}=\sgn_{S_n}$, the sign character of $S_n$. If $\alpha$ is a (finite) sequence of integers but is not a partition, then we interpret $\chi^\alpha$ to be the zero function.

The Young diagram of a partition $\lambda$ will be denoted by $[\lambda]$, and that of a skew shape $\lambda/\mu$ by $[\lambda/\mu]\coloneqq [\lambda]\setminus[\mu]$ for $\mu$ a subpartition of $\lambda$ (written $\mu\subseteq\lambda$). The boxes in a Young diagram will sometimes be referred to as nodes, and we refer to skew shapes and skew diagrams interchangeably when the meaning is clear.
We denote the length of the partition $\lambda$ by $l(\lambda)$, and the conjugate partition of $\lambda$ by $\lambda'$. Note
\begin{equation}\label{eqn:sign}
	\chi^{\lambda'}=\sgn_{S_n}\cdot\ \chi^\lambda
\end{equation}
for all $\lambda\vdash n$ (see \cite[2.1.8]{JK}). 

We record some operations on partitions. Let $\lambda=(\lambda_1,\lambda_2,\dotsc)$ and $\mu=(\mu_1,\mu_2,\dotsc)$ be two partitions. Then $+$ denotes component-wise addition, i.e.~$\lambda+\mu=(\lambda_1+\mu_1,\lambda_2+\mu_2,\dotsc)$, and $\lambda\sqcup\mu$ denotes the partition obtained by taking the disjoint union of the parts of $\lambda$ and $\mu$ and reordering so that parts are in non-increasing order.
When clear from context, we abbreviate $(a^b)\coloneqq (a,a,\dotsc,a)$ where there are $b$ parts of size $a$; in general we will specify $(a^b)\vdash ab$ to avoid confusion with the single part partition of $a^b$.

We also make use of skew characters for symmetric groups, i.e.~those indexed by skew shapes. For partitions $\mu$ and $\lambda$ such that $|\mu|\le|\lambda|$, the skew character $\chi^{\lambda/\mu}$ of $S_{|\lambda|-|\mu|}$ satisfies
\begin{equation}\label{eqn:skew}
	\langle \chi^{\lambda/\mu},\chi^\nu\rangle = \langle \chi^\lambda, (\chi^\mu\times \chi^\nu)\up_{S_{|\mu|}\times S_{|\nu|}}^{S_{|\mu|+|\nu|}}\rangle\qquad\forall\ \nu\vdash |\lambda|-|\mu|.
\end{equation}
Note if $\mu\not\subseteq\lambda$ then $\chi^{\lambda/\mu}=0$. This can be seen from the Littlewood--Richardson rule, which gives an explicit combinatorial description of the decomposition into irreducibles of the induced character appearing in the above expression, with Littlewood--Richardson coefficients arising as the multiplicities. These appear in many contexts, so we shall now fix the notation which will be used throughout this article (see \cite{J}).

\begin{definition}
	Let $n\in\N$, $\lambda=(\lambda_1,\dotsc,\lambda_k)\vdash n$ and $\mathcal{C}=(c_1,c_2,\dotsc,c_n)$ be a sequence of positive integers. We say that $\mathcal{C}$ is of type $\lambda$ if 
	\[ |\{i\in\{1,2,\dotsc,n\} \mid c_i=j\}|=\lambda_j\]
	for all $j\in\{1,2,\dotsc,k\}$. Moreover, we say that an element $c_j$ of $\mathcal{C}$ is good if $c_j=1$ or if
	\[ |\{i\in\{1,2,\dotsc,j-1\}\mid c_i=c_j-1\}| > |\{i\in\{1,2,\dotsc,j-1\}\mid c_i=c_j\}|. \]
	Finally, we say that $\mathcal{C}$ is good if $c_j$ is good for all $j\in\{1,\dotsc,n\}$.
\end{definition}

\begin{theorem}[Littlewood--Richardson rule]\label{thm:LR}
	Let $m,n\in\mathbb{N}_0$. Let $\mu\vdash m$ and $\nu\vdash n$. Then
	\[ (\chi^\mu\times\chi^\nu)\up^{S_{m+n}}_{S_m\times S_n} = \sum_{\lambda\vdash m+n} c^\lambda_{\mu,\nu}\cdot \chi^\lambda\]
	where the Littlewood--Richardson coefficient $c^\lambda_{\mu,\nu}$ equals the number of ways to replace the nodes of the skew Young diagram of $\lambda/\mu$ by natural numbers such that
	\begin{enumerate}[label=(\roman*)]
		\item\label{it1} the sequence obtained by reading the numbers from right to left, top to bottom is good of type $\nu$;
		\item \label{it2} the numbers are non-decreasing \textup{(}weakly increasing\textup{)} left to right along rows; and
		\item \label{it3} the numbers are strictly increasing down columns.
	\end{enumerate}
\end{theorem}

We call the order in Theorem~\ref{thm:LR}\ref{it1} the reading order of a skew shape. Let $\nu$ be a partition and $\gamma$ be a skew shape of size $|\nu|$. We call a way of replacing the nodes of $\gamma$ by numbers satisfying conditions Theorem~\ref{thm:LR}\ref{it1}--\ref{it3} a Littlewood--Richardson (LR) filling of $\gamma$ of type $\nu$. 
Clearly $c^\lambda_{\mu,\nu}=c^\lambda_{\nu,\mu}$. Using Littlewood--Richardson coefficients, we can also rephrase \eqref{eqn:skew} as $\chi^{\lambda/\mu}=\sum_\nu c^\lambda_{\mu,\nu}\cdot \chi^\nu$. Moreover, we can extend this notation to generalised Littlewood--Richardson coefficients $c^\lambda_{\mu^1,\dotsc,\mu^r}$ describing the constituents of
\[ (\chi^{\mu^1}\times\cdots\times\chi^{\mu^r})\up_{S_{n_1}\times\cdots\times S_{n_r}}^{S_{n_1+\cdots+n_r}}, \]
for any $r\in\N$ and $n_i\in\N_0$, and partitions $\mu^i\vdash n_i$ and $\lambda\vdash n_1+\cdots+n_r$. Similarly, $c^\lambda_{\mu^1,\dotsc,\mu^r}=c^\lambda_{\mu^{\sigma(1)},\dotsc,\mu^{\sigma(r)}}$ for any $\sigma\in S_r$. Furthermore, for $A\subseteq\cP(n)$ and $B\subseteq\cP(m)$ we define the operation $\star$ as follows:
\[ A\star B\coloneqq  \{\lambda\vdash m+n \mid \exists\ \mu\in A,\ \nu\in B\text{ s.t. }c^\lambda_{\mu,\nu}>0 \}. \]
We note that $\star$ is both commutative and associative.

Finally, for a partition $\lambda=(\lambda_1,\lambda_2,\dotsc,\lambda_k)\vdash n$, we let $S_\lambda\cong S_{\lambda_1}\times\cdots\times S_{\lambda_k}$ denote the corresponding Young subgroup of $S_n$. The permutation module $\triv_{S_\lambda}\up^{S_n}$ induced by the action of $S_n$ on the cosets of $S_\lambda$ will be denoted by $M^\lambda$. 
Young's Rule (see \cite[2.8.5]{JK}) tells us the decomposition of these permutation modules into irreducibles. Denoting the character of $M^\lambda$ by $\xi^\lambda$, we have that $\langle \xi^\lambda,\chi^\alpha\rangle$ equals the number of semistandard Young tableaux of shape $\alpha$ and content $\lambda$, for any $\alpha\vdash n$. Moreover, this multiplicity is positive if and only if $\alpha\trianglerighteq\lambda$, where $\trianglerighteq$ denotes the dominance partial order on partitions.



\subsection{Wreath products and Sylow subgroups of symmetric groups}
In order to describe the Sylow subgroups of symmetric groups, we briefly introduce some notation for wreath products. Let $G$ be a finite group and let $H\le S_n$ for some $n\in\N$. The natural action of $S_n$ on the factors of the direct product $G^{\times n}$ induces an action of $S_n$ (and therefore of $H$) via automorphisms of $G^{\times n}$, giving the wreath product $G\wr H\coloneqq  G^{\times n}\rtimes H$.  
As in \cite[Chapter 4]{JK}, we denote the elements of $G\wr H$ by $(g_1,\dotsc,g_n;h)$ for $g_i\in G$ and $h\in H$. Let $V$ be a $\bC G$--module and suppose it affords the character $\phi$. Let $V^{\otimes n}$ be the corresponding $\bC G^{\times n}$--module. The left action of $G\wr H$ on $V^{\otimes n}$ defined by linearly extending
\[ (g_1,\dotsc,g_n;h)\ :\quad v_1\otimes \cdots\otimes v_n \longmapsto g_1v_{h^{-1}(1)}\otimes\cdots\otimes g_nv_{h^{-1}(n)} \]
turns $V^{\otimes n}$ into a $\bC(G\wr H)$--module, which we denote by $\widetilde{V^{\otimes n}}$ (see \cite[(4.3.7)]{JK}), and we denote its character by $\tilde{\phi}$. For any $\psi\in\Ch(H)$, we define $\cX(\phi;\psi)$ as follows:
\[ \cX(\phi;\psi)\coloneqq \tilde{\phi}\cdot\Infl_H^{G\wr H}(\psi)\ \in\Ch(G\wr H). \]
The inflation $\Infl_H^{G\wr H}(\psi)$ of $\psi$ from $H$ to $G\wr H$ (identifying $H$ with the quotient $(G\wr H)/G^{\times n}$) is sometimes abbreviated to simply $\psi$, for convenience.

\begin{lemma}\label{lem:infl-ind}
	Let $G$, $H$ and $\phi$ be as above. Let $L\le H$ and $\tau\in\Ch(L)$. Then $\cX(\phi;\tau)\up_{G\wr L}^{G\wr H} = \cX(\phi;\tau\up_L^H)$.
\end{lemma}
\begin{proof}
	For any $\alpha\in\Ch(H)$ and $\beta\in\Ch(L)$, it is easy to check that $\alpha\cdot(\beta\up^H)=(\alpha\down_L\cdot \beta)\up^H$. Hence
	\[ \cX(\phi;\tau)\up_{G\wr L}^{G\wr H} \coloneqq  \big( \tilde{\phi}\down_{G\wr L}^{G\wr H}\cdot\Infl_L^{G\wr L}\tau\big)\up_{G\wr L}^{G\wr H} = \tilde{\phi} \cdot \big(\Infl_L^{G\wr L}(\tau)\big)\up_{G\wr L}^{G\wr H}. \]
	But induction and inflation of characters commute, so $\big(\Infl_L^{G\wr L}(\tau)\big)\up_{G\wr L}^{G\wr H} = \Infl_H^{G\wr H}(\tau\up_L^H)$. Thus $\cX(\phi;\tau)\up_{G\wr L}^{G\wr H} = \tilde{\phi}\cdot \Infl_H^{G\wr H}(\tau\up_L^H) = \cX(\phi;\tau\up_L^H)$, as claimed.
\end{proof}

Furthermore, if $\phi\in\Irr(G)$ then Gallagher's Theorem \cite[Corollary 6.17]{I} gives $\Irr(G\wr H\mid \phi^{\times n}) = \{\cX(\phi;\psi)\mid \psi\in\Irr(H)\}$, where $\Irr(G\wr H\mid \phi^{\times n})\coloneqq \{\chi\in\Irr(G\wr H)\mid \langle \chi\down_{G^{\times n}}, \phi^{\times n}\rangle \ne 0 \}$.
For a full description of $\Irr(G\wr H)$, we refer the reader to \cite[Chapter 4]{JK}; in the case $H=S_2$, we use the notation below.

\begin{notation}\label{not:GwrS2}
	Let $G$ be a finite group and suppose $\Irr(G)=\{\chi^i \mid i\in I\}$. Then 
	\[ \Irr(G\wr S_2) = \{ \psi^{i,j} \mid i\ne j\in I \} \sqcup \{ \psi^{i,i}_+, \psi^{i,i}_- \mid i\in I \} \]
	where 
	\[ \psi^{i,j} \coloneqq  (\chi^i\times\chi^j)\up_{G\times G}^{G\wr S_2}\ = \psi^{j,i},\quad \psi^{i,i}_{+} \coloneqq  \cX(\chi^i;\triv_{S_2})\quad\text{and}\quad \psi^{i,i}_{-} \coloneqq  \cX(\chi^i;\sgn_{S_2}). \]
\end{notation}

It will be useful to describe the decomposition of the permutation character $\triv_{H\wr S_2}\up^{G\wr S_2}$, for finite groups $H\le G$.

\begin{lemma}\label{lem:Prop1.2}
	Let $G$ be a finite group and let $\Irr(G)$ and $\Irr(G\wr S_2)$ be as in Notation~\ref{not:GwrS2}.
	Let $H\leq G$ and let $\pi \coloneqq  \triv_{H}\up^{G}$. Then $\tilde{\pi} \coloneqq  \triv_{H\wr S_2}\up^{G\wr S_2}$ decomposes into irreducible constituents with multiplicities given by
	\[ \langle\tilde{\pi},\psi^{i,j}\rangle = \langle\pi,\chi^i\rangle \cdot \langle\pi,\chi^j\rangle \quad\text{and}\quad \langle\tilde{\pi},\psi^{i,i}_{\pm}\rangle = \tfrac{1}{2}\cdot\big( \langle\pi,\chi^i\rangle^2 \pm \langle\pi,\chi^i\rangle \big). \]
\end{lemma}

\begin{proof}
	The first part follows from Mackey's theorem applied to the subgroups $G\times G$ and $H\wr S_2$ of $G\wr S_2$: since $(G\times G)\cdot (H\wr S_2)=G\wr S_2$ and $(G\times G)\cap (H\wr S_2)=H\times H$, we have
	\[ \langle\tilde{\pi},\psi^{i,j}\rangle	= \langle \triv_{H\wr S_2}\up^{G\wr S_2}\down_{G\times G}, \chi^i\times\chi^j\rangle = \langle \triv_{H\times H}\up^{G\times G}, \chi^i\times \chi^j\rangle = \langle\pi,\chi^i\rangle \cdot \langle\pi,\chi^j\rangle. \]
	Next, we use the wreath product character formula \cite[Lemma 4.3.9]{JK} to obtain
	\[ \psi^{i,i}_{\pm}(g_1,g_2;1) = \chi^i(g_1)\cdot\chi^i(g_2)\quad\text{and}\quad\psi^{i,i}_{\pm}\big(g_1,g_2;(1,2)\big) = \pm\chi^i(g_1 g_2) \qquad\forall\ g_1, g_2\in G. \]
	Hence
	\begin{align*}
		\langle\tilde{\pi},\psi^{i,i}_{\pm}\rangle &= \langle\triv_{H\wr S_2},\psi^{i,i}_{\pm}\down_{H\wr S_2}\rangle\\
		&= \frac{1}{|H\wr S_2|}\sum_{h_1,h_2\in H}\big(\chi^i(h_1)\cdot\chi^i(h_2)\pm\chi^i(h_1 h_2)\big) = \tfrac{1}{2}\big(\langle\pi,\chi^i\rangle^2\pm \langle\pi,\chi^i\rangle\big)
	\end{align*}
	as claimed.
\end{proof}

To describe Sylow subgroups of symmetric groups, fix a prime $p$ and let $n\in\N$. Let $P_n$ denote a Sylow $p$-subgroup of $S_n$. Clearly $P_1$ is trivial while $P_p$ is cyclic of order $p$. More generally, $P_{p^k}= (P_{p^{k-1}})^{\times p}\rtimes P_p=P_{p^{k-1}}\wr P_p\cong P_p\wr \cdots \wr P_p$ ($k$-fold wreath product) for all $k\in\N$. Now suppose $n\in\N$ and let $n=\sum_{i=1}^k a_ip^{n_i}$ be its $p$-adic expansion, i.e.~$n_1>\cdots>n_k\ge 0$ and $a_i\in\{1,2,\dotsc,p-1\}$ for all $i$. Then $P_n\cong (P_{p^{n_1}})^{\times a_1}\times\cdots\times (P_{p^{n_k}})^{\times a_k}$. 

To fix a convention for denoting such wreath products involving Sylow subgroups of symmetric groups more generally, we have the following.

\begin{notation}\label{not:convention}
	Let $G$ be a finite group and $p$ a prime. We use the convention that $P_n$ will always be viewed as a subgroup of $S_n$ in the notation $G\wr P_n$, that is, $G\wr P_n$ is a semi-direct product $G^{\times n}\rtimes P_n$. 
\end{notation}

\begin{remark}\label{rem:convention}
	Suppose $p=2$ and $n=2^{n_1}+\cdots+2^{n_k}$ for some $n_1>\cdots>n_k\ge 0$. With the convention of Notation~\ref{not:convention}, we observe that 
	\begin{align*}
		P_{2n} \cong P_{2^{n_1+1}}\times\cdots\times P_{2^{n_k+1}} &\cong (P_2\wr P_{2^{n_1}})\times\cdots\times (P_2\wr P_{2^{n_k}})\\
		&\cong P_2 \wr (P_{2^{n_1}}\times \cdots \times P_{2^{n_k}})\cong P_2 \wr P_n,
	\end{align*}
	viewing $P_{2^{n_1}}\times \cdots \times P_{2^{n_k}}\cong P_n$ naturally as a subgroup of $S_n$. Inductively, we also have $P_{2^tn}\cong P_{2^t}\wr P_n$ for all $t\in\N$.
	On the other hand, we clarify for example that $P_2\wr P_3 \not\cong P_2\wr P_2$ in this notation, even though $P_3\cong P_2$.%\hfill$\lozenge$
\end{remark}

We now return to an arbitrary prime $p$. Following the notation introduced in \cite{GL2}, given $\chi\in\Irr(S_n)$ and $\phi\in\Irr(P_n)$, the \textit{Sylow branching coefficient} $Z^\chi_\phi$ denotes the non-negative integer
\[ Z^\chi_\phi\coloneqq \langle \chi\down_{P_n}, \phi\rangle. \]
In this article, we will be particularly interested in the case where $\phi=\triv_{P_n}$, and abbreviate $Z^\chi_{\triv_{P_n}}$ to $Z^\chi$. Moreover, if $\chi=\chi^\lambda$ for a partition $\lambda$, then we shorten $Z^{\chi^\lambda}$ to $Z^\lambda$. 

We record one more lemma which will be useful later.
\begin{lemma}\label{lem:X-induced}
	Let $A$ and $B$ be finite groups, and let $n\in\N$. Then
	\[ \triv\up_{(A\times B)\wr S_n}^{A\wr S_n \times B\wr S_n} = \sum_{\phi\in\Irr(S_n)} \cX(\triv_A;\phi)\cdot\cX(\triv_B;\phi). \]
\end{lemma}
\begin{proof}
	Let $\phi\in\Irr(S_n)$. By Frobenius reciprocity, 
	\begin{align*}
		\langle \triv\up_{(A\times B)\wr S_n}^{A\wr S_n \times B\wr S_n}, \cX(\triv_A;\phi)\cdot\cX(\triv_B;\phi)\rangle &= \langle \triv, \big( \cX(\triv_A;\phi)\cdot\cX(\triv_B;\phi) \big) \down_{(A\times B)\wr S_n}\rangle\\
		&= \tfrac{1}{|(A\times B)\wr S_n|}\sum \cX(\triv_A;\phi)(x)\cdot\cX(\triv_B;\phi)(x),
	\end{align*}
where the sum runs over all $x\in (A\times B)\wr S_n$.
	But this equals $\tfrac{|A|^n\cdot|B|^n}{|(A\times B)\wr S_n|}\sum_{g\in S_n} \phi(g)^2$ by \cite[Lemma 4.3.9]{JK}. As symmetric group characters are real-valued (in fact, integer-valued), this simplifies to $\tfrac{1}{|S_n|}\sum_{g\in S_n}\overline{\phi(g)}\cdot\phi(g)=\langle \phi,\phi\rangle=1$.
\end{proof}



\subsection{Plethysms and deflations}\label{sec:plethysm-deflation}
When $\phi$ and $\psi$ are characters of symmetric groups, the characters $\cX(\phi;\psi)$ introduced above are closely related to plethysms of Schur functions: we give a brief description here and refer the reader to \cite{dBPW, Mac, Stanley} for further detail. Let $s_\lambda$ denote the Schur function corresponding to the partition $\lambda$, and $\circ$ the plethystic product of symmetric functions. Using the characteristic map (see e.g.~\cite[Chapter 7]{Stanley}) between class functions of finite symmetric groups and the ring of symmetric functions, we have the correspondence
\[ \cX(\chi^\nu;\chi^\lambda)\up_{S_{|\nu|}\wr S_{|\lambda|}}^{S_{|\nu|\cdot|\lambda|}}\ \longleftrightarrow\ s_\lambda\circ s_\nu \]
for all partitions $\lambda$ and $\nu$. Therefore the plethysm coefficient $a^\mu_{\lambda,\nu}$ satisfies
\begin{equation}\label{eqn:a}
	a^\mu_{\lambda,\nu} = \langle s_\lambda\circ s_\nu, s_\mu\rangle = \langle \cX(\chi^\nu;\chi^\lambda)\up_{S_{|\nu|}\wr S_{|\lambda|}}^{S_{|\nu|\cdot|\lambda|}}, \chi^\mu \rangle	
\end{equation}
for all partitions $\mu$ (and note that this is zero if $|\mu|\ne|\nu|\cdot|\lambda|$). 
We also introduce plethysm coefficients indexed by skew shapes: for partitions $\beta\subseteq\alpha$, $\delta\subseteq\gamma$ and $\nu$,
\begin{equation}\label{eqn:skew-plethysm}
	a^{\alpha/\beta}_{\gamma/\delta,\nu}\coloneqq  \sum_{\substack{\eta\vdash|\gamma|-|\delta|\\\zeta\vdash|\alpha|-|\beta|}} c^\gamma_{\eta,\delta}\cdot c^\alpha_{\zeta,\beta}\cdot a^\zeta_{\eta,\nu}.
\end{equation}
In other words, if $\phi$ and $\theta$ are skew shapes and $\nu$ is any partition, then $a^\phi_{\theta,\nu} = \langle \cX(\chi^\nu; \chi^\theta)\up_{S_{|\nu|}\wr S_{|\theta|}}^{S_{|\nu|\cdot|\theta|}}, \chi^\phi\rangle$, extending the equality in \eqref{eqn:a}. We also remark that if $\zeta$ and $\eta$ are partitions, then $a^\zeta_{\eta,\emptyset}=1$ if $\zeta=\emptyset$ and $\eta=(n)$ for some $n$, and $a^\zeta_{\eta,\emptyset}=0$ otherwise. 

A well known symmetry property of plethysm coefficients involving the conjugation involution is the following (see e.g.~\cite[Ex.~1, Ch. I.8]{Mac}):
\begin{lemma}\label{lem:conj-plethysm}
	Let $\lambda$, $\mu$ and $\nu$ be partitions. Then
	\[ a^\nu_{\lambda,\mu} = a^{\nu'}_{\lambda^\ast,\mu'},\quad\text{where}\quad \lambda^\ast\coloneqq \begin{cases} \lambda & \text{if $|\mu|$ is even},\\ \lambda' & \text{if $|\mu|$ is odd}.\end{cases} \]
\end{lemma}

Character deflations were introduced in \cite[Definition 1.1]{EPW}, and used to prove results generalising the Murnaghan--Nakayama rule for computing symmetric group character values and (special cases of) the Littlewood--Richardson rule, as well as to verify new cases of the long-standing Foulkes' Conjecture. We observe that they give another language in which to describe certain plethysm coefficients. We first record the definition of these deflations in notation which we have introduced thus far.
\begin{definition}\label{def:deflation}
	Let $m,n\in\N$ and $\theta\in\Irr(S_m)$. Let $\xi\in\Irr(S_m\wr S_n)$. Then 
	\[ \operatorname{Def}^\theta_{S_n}(\xi) \coloneqq  \begin{cases}
		\chi^\nu & \text{ if }\xi = \cX(\theta;\chi^\nu)\text{ for some }\nu\vdash n,\\
		0 & \text{ otherwise},
	\end{cases} \]
	which then extends linearly to all class functions of $S_m\wr S_n$. If $\chi$ is a class function of $S_{mn}$ then set
	\[ \operatorname{Defres}^\theta_{S_n}(\chi) \coloneqq  \operatorname{Def}^\theta_{S_n}(\chi\down_{S_m\wr S_n}). \]
	When $n\in\N$ is fixed and $\gamma\vdash mn$ for some $m\in\N$, we use the notation
	\[ \delta^\gamma\coloneqq \operatorname{Defres}^{\triv_{S_m}}_{S_n}(\chi^\gamma) = \operatorname{Def}^{\triv_{S_m}}_{S_n}(\chi^\gamma\down^{S_{mn}}_{S_m\wr S_n}). \]
	In this article, we will sometimes refer to $\delta^\gamma$ as the deflation of $\gamma$ with respect to $S_n$ \textup{(}where the $m$ is understood from $|\gamma|/n$ and we suppress $m$ from the notation\textup{)}.
\end{definition}
In particular, if $\lambda\vdash n$ then we have that $a^\gamma_{\lambda,(m)} = \langle \delta^\gamma,\chi^\lambda \rangle$. This relation between plethysm coefficients and deflations can immediately be extended to skew shapes using Littlewood--Richardson coefficients. Namely, if $\alpha,\beta,\nu,\lambda$ are partitions such that $|\alpha|-|\beta|=n$ and $|\nu|-|\lambda|=mn$, then we may define $\delta^{\nu/\lambda}\coloneqq \operatorname{Defres}^{\triv_{S_m}}_{S_n}(\chi^{\nu/\lambda}) = \operatorname{Def}^{\triv_{S_m}}_{S_n}(\chi^{\nu/\lambda}\down^{S_{mn}}_{S_m\wr S_n})$. Then $\chi^{\alpha/\beta}=\sum_\sigma c^\alpha_{\beta,\sigma}\cdot \chi^\sigma$ and $\delta^{\nu/\lambda}=\sum_\tau c^\nu_{\lambda,\tau}\cdot \delta^\tau$ give
\begin{equation}\label{eqn:plethysm-deflation-skew}
	a^{\nu/\lambda}_{\alpha/\beta,(m)} = \langle \delta^{\nu/\lambda}, \chi^{\alpha/\beta}\rangle.
\end{equation}

The following is a straightforward result on plethysm coefficients involving `tall' partitions. We include a proof for convenience.
\begin{lemma}\label{lem:tall-pleth}
	Let $m,n\in\N$, $\lambda\vdash mn$ and $\nu\vdash n$. If $l(\lambda)>n$, then $a^\lambda_{\nu,(m)}=0$.
\end{lemma}

\begin{proof}
	Let $a\coloneqq a^\lambda_{\nu,(m)}= \langle \chi^\lambda\down^{S_{mn}}_{S_m\wr S_n}, \cX(\chi^{(m)};\chi^\nu)\rangle$. Let $Y\coloneqq (S_m)^{\times n}\le S_m\wr S_n$. 
	Let $\cX\coloneqq \cX(\chi^{(m)};\chi^\nu)$ and write $\chi^\lambda\down_{S_m\wr S_n} = a\cX + \Delta$ for some $\Delta\in\Ch(S_m\wr S_n)$. Since $\cX\down_Y = \chi^\nu(1)\cdot(\chi^{(m)})^n=\chi^\nu(1)\cdot\triv_Y$, it follows that $\langle \chi^\lambda\down_Y,\triv_Y\rangle \ge \langle a\cdot \cX\down_Y, \triv_Y\rangle = a\cdot\chi^\nu(1).$
	But by the Littlewood--Richardson rule, $l(\lambda)>n$ implies $\langle \chi^\lambda\down_Y,\triv_Y\rangle=0$ since $\triv_Y = (\chi^{(m)})^n$, which gives $a=0$.
\end{proof}

Finally, we record the following result of Thrall \cite{T}. We note that a partition is even if all of its parts are of even size.
\begin{proposition}\label{prop:thrall}
	Let $n\in\N$. Then 
	\begin{enumerate}[label=\textup{(\roman*)}]
		\item $s_{(n)}\circ s_{(2)} = \sum_\lambda s_\lambda$ where the sum runs over all even partitions $\lambda\vdash 2n$, and
		\item $s_{(2)}\circ s_{(n)} = \sum_\mu s_\mu$ where the sum runs over all even partitions $\mu\vdash 2n$ with at most two parts.
	\end{enumerate}
\end{proposition}
In other words, $a^\lambda_{(n),(2)}=1$ if $\lambda\vdash 2n$ is even and $a^\lambda_{(n),(2)}=0$ otherwise, while $a^\mu_{(2),(n)}=1$ if $\mu\vdash 2n$ is even and $l(\mu)\le 2$, and $a^\mu_{(2),(n)}=0$ otherwise.



% END NATIVE AUTHOR BLOCK
\setcounter{section}{3}
% BEGIN NATIVE AUTHOR BLOCK 628--676
\section{Plethysms and character deflations}\label{sec:pre-recursive-formula}
We record a result of Briand--Orellana--Rosas on plethysm coefficients from which we can deduce the $k=0$ case of Theorem A (Theorem~\ref{thm:4.4.3}), as well as resolve two conjectures of de Boeck in Section~\textup{4.1} (Theorems~\textup{4.5} and~\textup{4.6}).

\begin{definition}
	Let $\lambda$ be a partition and $w,h\in\N_0$ such that $\lambda\subseteq (w^h)$. Then $\square_{w,h}(\lambda)$ denotes the partition of size $wh-|\lambda|$ whose Young diagram is the $180^\circ$ rotation of the complement of the Young diagram of $\lambda$ in a rectangle of width $w$ and height $h$.
\end{definition}

\begin{theorem}[{\cite[Theorem 12]{BOR}}]\label{thm:BOR}
	Let $m,n\in\N_0$ and let $\lambda$, $\mu$ and $\nu$ be partitions such that $\mu\subseteq(w^h)$ and $l(\nu)\le h$. If $\nu\subseteq \big((w|\lambda|)^h\big)$ then
	\[ a^\nu_{\lambda,\mu} = a^{\square_{w|\lambda|,h}(\nu)}_{\lambda,\square_{w,h}(\mu)}. \]
	Otherwise $a^\nu_{\lambda,\mu}=0$.
\end{theorem}

Let $m,n\in\N$ and $\mu\vdash mn$. We recall the definition of $\delta^\mu=\operatorname{Def}^{\triv_{S_m}}_{S_n}(\chi^\mu\down^{S_{mn}}_{S_m\wr S_n})$, the deflation of $\mu$ with respect to $S_n$, from Definition~\ref{def:deflation}.
\begin{proposition}\label{prop:12.4}
	Fix $n\in\N$ and $m_1,m_2\in\N$. Let $\mu_1\vdash m_1n$ and suppose $\mu_1\subseteq \big((m_1+m_2)^n\big)$. Set $\mu_2=\square_{m_1+m_2,n}(\mu_1)\vdash m_2n$. Then
	\[ \delta^{\mu_1} = \begin{cases}
		\delta^{\mu_2} & \text{if }m_1+m_2\text{ is even},\\
		\sgn_{S_n}\cdot\ \delta^{\mu_2} & \text{if }m_1+m_2\text{ is odd},
	\end{cases}
	\]
	where $\delta^{\mu_i}$ refers to the deflation of $\mu_i$ with respect to $S_n$, for $i\in\{1,2\}$.
\end{proposition}

\begin{proof}
	Since $(1^{m_1})\subseteq(1^{m_1+m_2})$ and $l(\mu'_1)\le m_1+m_2$, as well as $\mu'_1\subseteq \big((1\cdot|\lambda|)^{m_1+m_2}\big)$, we see from Theorem~\ref{thm:BOR} that
	\[ a^{\mu_1'}_{\lambda,(1^{m_1})} = a^{\square_{n,m_1+m_2}(\mu_1')}_{\lambda,\square_{1,m_1+m_2}\big((1^{m_1})\big)} = a^{\mu'_2}_{\lambda,(1^{m_2})}\qquad\forall\ \lambda\vdash n. \]
	Applying Lemma~\ref{lem:conj-plethysm} and observing that $\lambda\mapsto\lambda'$ is an involution on the set of partitions of $n$, we obtain
	\[ a^{\mu_1}_{\lambda,(m_1)} = a^{\mu_2}_{\lambda^\ast,(m_2)} \]
	for all $\lambda\vdash n$, where $\lambda^\ast$ denotes $\lambda$ (resp.~$\lambda'$) if $m_1+m_2$ is even (resp.~odd). 
	Since $\langle \delta^\mu, \chi^\lambda\rangle = a^\mu_{\lambda,(m)}$, we have shown that
	\[ \langle \delta^{\mu_1},\chi^\lambda\rangle = \begin{cases}
		\langle \delta^{\mu_2},\chi^\lambda\rangle & \text{if $m_1+m_2$ is even},\\ \langle \delta^{\mu_2}, \chi^{\lambda'}\rangle & \text{if $m_1+m_2$ is odd}.\end{cases} \]
	Since $\chi^{\lambda'}=\sgn_{S_n}\cdot\ \chi^\lambda$, we have $\delta^{\mu_1}=\delta^{\mu_2}$ if $m_1+m_2$ is even (resp.~$\delta^{\mu_1}=\sgn_{S_n}\cdot\ \delta^{\mu_2}$ if $m_1+m_2$ is odd), as required.
\end{proof}

The following result was first proven in \cite[Proposition 1.16]{BCV}, and is a special case of \cite[Theorem 1.1]{dBPW}. It also forms the $k=0$ case of Theorem A, since $\sum_{\alpha,\beta\vdash 0} \big(a^{\alpha/\emptyset}_{\beta',(m)}\cdot a^{\hat{\mu}/\alpha}_{\lambda/\beta,(m-1)}\big) = a^{\hat{\mu}}_{\lambda,(m-1)}$. We give a short proof in the language of deflations.

\begin{theorem}\label{thm:4.4.3}
	Let $m,n\in\N$, $\lambda\vdash mn$ and $\nu\vdash n$. Then $a^\lambda_{\nu,(m)} = a^{\lambda+(1^n)}_{\nu',(m+1)}$.
\end{theorem}

\begin{proof}
	From Lemma~\ref{lem:tall-pleth}, $a^\lambda_{\nu,(m)} = a^{\lambda+(1^n)}_{\nu',(m+1)}=0$ if $l(\lambda)>n$, so we may now assume $l(\lambda)\le n$. 
	Since $\chi^{\nu'}=\sgn_{S_n}\cdot\ \chi^\nu$, the assertion is equivalent to proving that $\delta^\lambda=\sgn_{S_n}\cdot\ \delta^{\lambda+(1^n)}$. 
	Choose $k\in\N$ such that $\lambda\subseteq(k^n)$, say $k=mn$. Then by Proposition~\ref{prop:12.4},
	\[ \delta^\lambda =s_1\cdot \ \delta^{\square_{k,n}(\lambda)} = s_1\cdot s_2\cdot\ \delta^{\square_{1+k,n}(\square_{k,n}(\lambda))} = s_1\cdot s_2\cdot\ \delta^{\lambda+(1^n)}, \]
	where $\{s_1,s_2\}=\{\triv_{S_n},\sgn_{S_n}\}$.
\end{proof}
% END NATIVE AUTHOR BLOCK
\setcounter{section}{4}
% BEGIN NATIVE AUTHOR BLOCK 754--1082
\section{A recursive formula for plethysm coefficients}\label{sec:14}
The main result of this article is Theorem A, a recursive formula for plethysm coefficients of the form $a^\mu_{\lambda,(m)}$ for arbitrary $m\in\N$ and partitions $\mu$ and $\lambda$. Together with Lemma~\ref{lem:tall-pleth}, it describes the deflations $\delta^\mu$ of $\mu\vdash mn$ with respect to $S_n$, noting that $a^\mu_{\lambda',(m)}=\langle \sgn_{S_n}\cdot\ \delta^\mu,\chi^\lambda\rangle$ for $\lambda\vdash n$. We restate Theorem A here as Theorem~\ref{thm:14.6i} for ease of reference for the reader, and recall that plethysm coefficients indexed by skew shapes were defined in \eqref{eqn:skew-plethysm}:

\begin{theorem}[Theorem A]\label{thm:14.6i}
	Fix $n\in\N$. Let $m\in\N$, $k\in\{0,1,\dotsc,n-1\}$ and $\lambda\vdash n$. Let $\mu\vdash mn$ with $l(\mu)=n-k$, and set $\hat{\mu}\coloneqq \mu-(1^{n-k})\vdash (m-1)n+k$. Then
	\begin{equation}\label{eqn:A}
		a^\mu_{\lambda',(m)} = \sum_{i=0}^k (-1)^{k+i} \cdot \sum_{\substack{\alpha\vdash k+(m-1)i\\\beta\vdash i}} a^{\alpha/(k-i)}_{\beta',(m)}\cdot a^{\hat{\mu}/\alpha}_{\lambda/\beta,(m-1)}.
	\end{equation}
\end{theorem}

We first illustrate some of the uses of our main theorem in Example~\ref{ex:14.7} and in proving a stability result (Proposition~\ref{prop:17}), before proving Theorem~\ref{thm:14.6i} in full in Section~\ref{sec:16}. We present further applications to Sylow branching coefficients in Section~\textup{6}.

\begin{example}\label{ex:14.7}
	We illustrate how to compute $a^{(6,3^{n-2})}_{\lambda,(3)}$ for all $n\ge 6$ and $\lambda\vdash n$ using Theorem~\ref{thm:14.6i}. 
	Let $\mu=(6,3^{n-2})$. Since $l(\mu)=n-1$, Theorem~\ref{thm:14.6i} gives
	\begin{align*} a^\mu_{\lambda,(3)}& = - a^{\hat{\mu}/(1)}_{\lambda',(2)} + a^{\hat{\mu}/(3)}_{\lambda'/(1),(2)} \\&= -a^{(4,2^{n-2})}_{\lambda',(2)} - a^{(5,2^{n-3},1)}_{\lambda',(2)} + a^{(2^{n-1})}_{\lambda'/(1),(2)} + a^{(3,2^{n-3},1)}_{\lambda'/(1),(2)} + a^{(4,2^{n-3})}_{\lambda'/(1),(2)}, \end{align*}
	since $a^\emptyset_{\emptyset,(3)}=1$ and $a^\alpha_{(1),(3)}=\delta_{\alpha,(3)}$. Applying Theorems~\ref{thm:14.6i} and~\ref{thm:LR}, we obtain
	\begin{align*}
		a^{(4,2^{n-2})}_{\lambda',(2)} &= \begin{cases} 1 & \text{if }\lambda'\in\{(n),(n-1,1),(n-2,2)\},\\ 0&\text{otherwise},\end{cases}\\
		a^{(5,2^{n-3},1)}_{\lambda',(2)} &= \begin{cases} 1 & \text{if }\lambda'\in\{(n-1,1),(n-2,2),(n-2,1^2),(n-3,2,1)\},\\ 0&\text{otherwise}.\end{cases}
	\end{align*}
	Similarly, since $a^\theta_{\lambda'/(1),(2)} = \sum_{\varepsilon\vdash n-1}c^{\lambda'}_{\varepsilon,(1)}\cdot a^\theta_{\varepsilon,(2)}$, we have that
	\begin{align*}
		a^{(2^{n-1})}_{\lambda'/(1),(2)} &= c^{\lambda'}_{(n-1),(1)} = \begin{cases}
			1 & \text{if }\lambda'\in\{(n),(n-1,1)\},\\ 0&\text{otherwise},\end{cases}\\
		a^{(3,2^{n-3},1)}_{\lambda'/(1),(2)} &= c^{\lambda'}_{(n-2,1),(1)} = \begin{cases}
			1 & \text{if }\lambda'\in\{(n-1,1),(n-2,2),(n-2,1^2)\},\\ 0&\text{otherwise},\end{cases}\\
		a^{(4,2^{n-3})}_{\lambda'/(1),(2)} &= c^{\lambda'}_{(n-2,1),(1)} + c^{\lambda'}_{(n-1),(1)} + c^{\lambda'}_{(n-3,2),(1)}\\
		&= \begin{cases}
			2 & \text{if }\lambda'\in\{(n-1,1),(n-2,2)\},\\
			1 & \text{if }\lambda'\in\{(n),(n-2,1^2),(n-3,3),(n-3,2,1)\},\\
			0 & \text{otherwise}.
		\end{cases}
	\end{align*}
	Putting this together, we obtain
	\[ a^\mu_{\lambda,(3)} = \begin{cases}
		2 & \text{if }\lambda'=(n-1,1),\\
		1 & \text{if }\lambda'\in\{(n),(n-2,2),(n-2,1^2),(n-3,3)\},\\
		0 & \text{otherwise}.
	\end{cases} \]
	In particular, this gives an alternative method for one of the steps in the proof of Theorem~\textup{4.6} by showing that $a^{(6,3^{n-2})}_{(1^n),(3)}=1$ for all $n\ge 3$ (note when $n\le 5$ this follows by direct computation).%\hfill$\lozenge$
\end{example}

A corollary of Theorem~\ref{thm:14.6i} is the stability of the following sequence of plethysm coefficients, whose monotonicity was predicted in \cite[Conjecture 1.2]{BBP}.

\begin{proposition}\label{prop:17}
	Let $\lambda$ and $\mu$ be partitions. For all $j\in\N_0$, define partitions $\lambda^j\coloneqq \lambda\sqcup(1^j)$ and $\mu^j\coloneqq (\mu+(j))\sqcup(1^j)$. Then the sequence $\big(a^{\mu^j}_{\lambda^j,(2)}\big)_{j\in\N_0}$ is eventually constant.
\end{proposition}

To prove Proposition~\ref{prop:17}, we first record a stability property of Littlewood--Richardson coefficients. For convenience we include a proof in our present notation.
\begin{lemma}\label{lem:lr-stab}
	Let $\alpha,\beta,\gamma$ and $\delta$ be partitions. Define $\gamma(j)\coloneqq \gamma+(j)$ and $\delta(j)\coloneqq \delta+(j)$ for all $j\in\N_0$. Then the sequence $\big( \langle\chi^{\gamma(j)/\alpha},\chi^{\delta(j)/\beta}\rangle \big)_{j\in\N_0}$ is non-decreasing and eventually constant.
\end{lemma}

\begin{proof}
	For a skew shape $\rho$, let $\cL(\rho)$ denote the set of all Littlewood--Richardson fillings of $\rho$. Let $\ff_j:\cL(\gamma(j)/\alpha)\to \cL(\gamma(j+1)/\alpha)$ be the map given by filling the additional box in $\gamma(j+1)/\alpha - \gamma(j)/\alpha$ with the number 1. Clearly $\ff_j$ is well-defined and injective for all $j$, and furthermore, bijective for all sufficiently large $j$ (for example, $j\ge \alpha_1+\gamma_2$ will suffice).
	Similarly define $\fg_j:\cL(\delta(j)/\beta)\to\cL(\delta(j+1)/\beta)$. For $\mathsf{s}\in\cL(\gamma(j)/\alpha)$ and $\mathsf{t}\in\cL(\delta(j)/\beta)$, note that $\mathsf{s}$ and $\mathsf{t}$ have the same type if and only if $\ff_j(\mathsf{s})$ and $\fg_j(\mathsf{t})$ have the same type. The assertion of the lemma follows since $\langle\chi^{\gamma(j)/\alpha},\chi^{\delta(j)/\beta}\rangle = |\{(\mathsf{s},\mathsf{t})\in\cL(\gamma(j)/\alpha)\times\cL(\delta(j)/\beta) \mid \mathsf{s}, \mathsf{t}\text{ have the same type} \}|$.
\end{proof}

\begin{proof}[Proof of Proposition~\ref{prop:17}]
	We may assume that $\lambda\vdash n$, $\mu\vdash 2n$ and $l(\mu)=n-k$ for some $n$ and $k\in\N_0$, else $a^{\mu^j}_{\lambda^j,(2)}=0$ for all $j$ by Lemma~\ref{lem:tall-pleth}. Since $\mu^j\vdash 2n+2j$ and $l(\mu^j)=(n+j)-k$, we have from Theorem~\ref{thm:14.6i} that
	\[ a^{\mu^j}_{\lambda^j,(2)} = \sum_{i=0}^k (-1)^{k+i}\cdot \sum_{\substack{\alpha\vdash k+i\\\beta\vdash i}} a^{\alpha/(k-i)}_{\beta',(2)}\cdot a^{\hat{\mu^j}/\alpha}_{{\lambda^j}'/\beta,(1)} \]
	for all $j$, where $\hat{\mu^j}\coloneqq \mu^j-(1^{n+j-k})$. The proof is concluded by observing that $a^{\hat{\mu^j}/\alpha}_{{\lambda^j}'/\beta,(1)}=\langle \chi^{{\lambda^j}'/\beta},\chi^{\hat{\mu^j}/\alpha}\rangle$ and using Lemma~\ref{lem:lr-stab} with $\gamma\coloneqq \hat{\mu}$ and $\delta\coloneqq \lambda'$.
\end{proof}

\begin{remark}
	A similar argument can be used to give a new proof that the sequence $\big(a^{\mu^j}_{\lambda^j,(2)}\big)_j$ also stabilises where $\lambda^j\coloneqq \lambda+(j)$ and $\mu^j=\mu\sqcup(2^j)$; this sequence is already known to be both non-decreasing and eventually constant by \cite[\textsection 2.6 Corollary 1]{Brion}.%\hfill$\lozenge$
\end{remark}



\subsection{Proof of Theorem~\ref{thm:14.6i}}\label{sec:16}
We first introduce some notation in preparation for the proofs to come.
\begin{notation}\label{not:rho}
	\begin{enumerate}[label=\textup{(\roman*)}]
		\item \label{561} Let $\lambda$ be a partition, $n\in\N_0$ and $\phi$ be a virtual character of $S_n$.
		We define
		\[ \phi/\chi^\lambda\coloneqq  \sum_{\mu\vdash n}\langle \phi,\chi^\mu\rangle\cdot\chi^{\mu/\lambda} \]
		where $\chi^{\mu/\lambda}=0$ if $\lambda\not\subseteq\mu$.
		
		\item \label{562} For $m,n\in\N$ and $\alpha/\beta$ a skew shape of size $n$, define \[
		\rho^{\alpha/\beta}_m \coloneqq  \cX(\triv_{S_m};\chi^{\alpha/\beta})\up_{S_m\wr S_n}^{S_{mn}}.\] 
		
		\item \label{563} For $\phi\in\Ch(S_{n_1})$ and $\theta\in\Ch(S_{n_2})$, define
		\[ \phi\boxtimes\theta\coloneqq  (\phi\times\theta)\up_{S_{n_1}\times S_{n_2}}^{S_{n_1+n_2}}. \]
		
		\item \label{564} Let $S_\lambda$ denote the Young subgroup $S_{\lambda_1}\times\cdots\times S_{\lambda_{l(\lambda)}}$ of $S_n$ and let
		\[ \zeta^\lambda\coloneqq \triv_{S_\lambda}\up^{S_n} = \chi^{(\lambda_1)}\boxtimes\cdots\boxtimes \chi^{(\lambda_{l(\lambda)})} \]
		denote the character of the permutation module $M^\lambda$. (Note $\zeta^\mu=\zeta^\lambda$ if $\mu$ is a composition of $n$ with the same parts as $\lambda$ but in a different order.) 
	\end{enumerate}
\end{notation}

Recall that the irreducible decomposition of such a permutation character $\zeta^\lambda$ is described by Young's Rule \cite[2.8.5]{JK}. Equivalently,
\begin{equation}\label{eqn:kostka}
	\zeta^\lambda = \sum_{\gamma\vdash n} K_{\gamma,\lambda}\cdot \chi^\gamma
\end{equation}
where the Kostka number $K_{\gamma,\lambda} = \langle \zeta^\lambda,\chi^\gamma\rangle=c^\gamma_{(\lambda_1),\dotsc,(\lambda_{l(\lambda)})}$ equals the number of semistandard Young tableaux of shape $\gamma$ and content $\lambda$.
In particular, we therefore have
\begin{equation}\label{eqn:15.2}
	\sum_{\gamma\vdash n} K_{\gamma,\lambda}\cdot \rho^\gamma_m = \cX(\triv_{S_m}; \zeta^\lambda)\up_{S_m\wr S_n}^{S_{mn}} = \rho^{(\lambda_1)}_m\boxtimes\cdots\boxtimes\rho^{(\lambda_{l(\lambda)})}_m.
\end{equation}

We now precisely identify the $i=k$ term of Theorem~\ref{thm:14.6i}, in Theorem~\ref{thm:14.6ii}. We then deduce Theorem~\ref{thm:14.6i} from Theorem~\ref{thm:14.6ii}. Following that, we prove Theorem~\ref{thm:14.6ii}, during the course of which we will also prove Theorem B, which has been numbered as Theorem~\ref{thm:16.7} in this section for ease of reference.

\begin{theorem}\label{thm:14.6ii}
	Fix $n\in\N$. Let $m\in\N$, $k\in\{0,1,\dotsc,n-1\}$ and $\lambda\vdash n$. Let $\nu\vdash mn+k$ with $l(\nu)=n$, and set $\hat{\nu}\coloneqq \nu-(1^n)\vdash (m-1)n+k$. Then
	\[ a^{\nu/(1^k)}_{\lambda',(m)} = \sum_{\substack{\alpha\vdash mk\\\beta\vdash k}} a^\alpha_{\beta',(m)}\cdot a^{\hat{\nu}/\alpha}_{\lambda/\beta,(m-1)}. \]
\end{theorem}

\begin{proof}[Proof of Theorem~\ref{thm:14.6i} from Theorem~\ref{thm:14.6ii}]
	We proceed by induction on $k$, observing that the $k=0$ case of Theorem~\ref{thm:14.6i} follows immediately from the $k=0$ case of Theorem~\ref{thm:14.6ii}. 
	Now assume $k>0$, and fix $\mu\vdash mn$ with $l(\mu)=n-k$. Let $\nu=\mu\sqcup(1^k)$, so $l(\nu)=n$ and $\hat{\nu}=\hat{\mu}$.
	
	We aim to evaluate $a^{\nu/(1^k)}_{\lambda',(m)}$, and compare it to Theorem~\ref{thm:14.6ii}. To do this, we will study the constituents in the skew character $\chi^{\nu/(1^k)}$. First, note that for any $\omega\vdash|\nu|-k$, by Theorem~\ref{thm:LR} we must have $c^\nu_{\omega,(1^k)}\in\{0,1\}$. Moreover, $c^\nu_{\omega,(1^k)}=1$ if and only if $\omega\subseteq\nu$ and all $k$ boxes of $[\nu/\omega]$ belong to different rows. We will denote by $\mathcal{A}$ the collection of $\omega$ with $c^\nu_{\omega,(1^k)}=1$. In particular,
	\[ \sum_{\omega\in\mathcal{A}}\chi^\omega = \chi^{{\nu}/(1^k)}. \]
	We partition $\mathcal{A}$ as a disjoint union, $\mathcal{A} = \coprod_{j=0}^{k}\mathcal{A}_j$, where $\mathcal{A}_j$ is the collection of $\omega\in\mathcal{A}$ for which $l(\omega)=n-k+j$. Notice that for each $j$, $\mathcal{A}_j$ bijects to $\mathcal{B}_j \coloneqq  \{\varpi\vdash|\hat{\mu}|-j \mid c^{\hat{\mu}}_{\varpi,(1^j)}\}$; the map is given by removal of the first column, $\omega\mapsto\hat{\omega}$, and this is seen to be a bijection by Theorem~\ref{thm:LR}. By a similar application of Theorem~\ref{thm:LR},
	
	\begin{equation}\label{eqn:B_j}
		\sum_{\varpi\in\mathcal{B}_j}\chi^\varpi = \chi^{\hat{\mu}/(1^j)}.
	\end{equation}
	Now observe that
	\[ a^{\nu/(1^k)}_{\lambda',(m)} = \sum_{\omega\in\mathcal{A}} a^{\omega}_{\lambda',(m)} = a^{\mu}_{\lambda',(m)} + \sum_{j=1}^k \sum_{\omega\in\mathcal{A}_j} a^{\omega}_{\lambda',(m)}. \]
	The idea here is that in our partition of $\mathcal{A}$, the $j=0$ term contributes precisely $a^{\mu}_{\lambda',(m)}$. Let us set $X \coloneqq  -\sum_{j=1}^k \sum_{\omega\in\mathcal{A}_j} a^{\omega}_{\lambda',(m)}$. Assuming Theorem~\ref{thm:14.6ii}, it suffices to show that $X$ equals the $\sum_{i=0}^{k-1}(\dotsc)$ part of the summation on the right hand side of \eqref{eqn:A}. For $0<j\leq k$, we have that
	\begin{align*}
		\sum_{\omega\in\mathcal{A}_j} a^{\omega}_{\lambda',(m)} &\overset{\text{\normalfont\tiny(inductive hypothesis)}}{=} \sum_{\omega\in\mathcal{A}_j} \sum_{i=0}^{k-j}(-1)^{k-j+i}\sum_{\substack{\gamma\vdash k-j + (m-1)i\\ \beta\vdash i}}a^{\gamma/(k-j-i)}_{\beta',(m)}\cdot a^{\hat{\omega}/\gamma}_{\lambda/\beta,(m-1)}\\
		&\overset{\text{\normalfont\tiny(bijection $\mathcal{A}_j\to\mathcal{B}_j$)}}{=}\sum_{\varpi\in\mathcal{B}_j} \sum_{i=0}^{k-j}(-1)^{k-j+i}\sum_{\substack{\gamma\vdash k-j + (m-1)i\\ \beta\vdash i}}a^{\gamma/(k-j-i)}_{\beta',(m)}\cdot a^{\varpi/\gamma}_{\lambda/\beta,(m-1)}
		\\
		&\overset{\normalfont\tiny \text{\eqref{eqn:a}}}{=}\sum_{\varpi\in\mathcal{B}_j} \sum_{i=0}^{k-j}(-1)^{k-j+i}\sum_{\substack{\gamma\vdash k-j + (m-1)i\\ \beta\vdash i}}a^{\gamma/(k-j-i)}_{\beta',(m)}\cdot \left\langle\chi^{\varpi},\rho^{\lambda/\beta}_{m-1}\boxtimes\chi^\gamma\right\rangle \\
		&\overset{\normalfont\tiny\text{\eqref{eqn:B_j}}}{=}\sum_{i=0}^{k-j}(-1)^{k-j+i}\sum_{\substack{\gamma\vdash k-j + (m-1)i\\ \beta\vdash i}}a^{\gamma/(k-j-i)}_{\beta',(m)}\cdot \left\langle\chi^{\hat{\mu}/(1^j)},\rho^{\lambda/\beta}_{m-1}\boxtimes\chi^\gamma\right\rangle
		\\
		&=\sum_{i=0}^{k-j}(-1)^{k-j+i}\sum_{\substack{\gamma\vdash k-j + (m-1)i\\ \beta\vdash i}} a^{\gamma/(k-j-i)}_{\beta',(m)}\cdot \left\langle\chi^{\hat{\mu}},\rho^{\lambda/\beta}_{m-1}\boxtimes\chi^\gamma\boxtimes\chi^{(1^j)}\right\rangle\\
		&=\sum_{i=0}^{k-j}(-1)^{k-j+i}\sum_{\substack{\gamma\vdash k-j + (m-1)i\\ \alpha\vdash k+(m-1)i\\ \beta\vdash i}} c^{\alpha}_{\gamma,(1^j)}\cdot a^{\gamma/(k-j-i)}_{\beta',(m)}\cdot \left\langle\chi^{\hat{\mu}},\rho^{\lambda/\beta}_{m-1}\boxtimes\chi^\alpha\right\rangle\\
		&\overset{\normalfont\tiny\text{\eqref{eqn:a}}}{=}\sum_{i=0}^{k-j}(-1)^{k-j+i}\sum_{\substack{\gamma\vdash k-j + (m-1)i\\ \alpha\vdash k+(m-1)i\\ \beta\vdash i}} c^{\alpha}_{\gamma,(1^j)}\cdot a^{\gamma/(k-j-i)}_{\beta',(m)}\cdot a^{\hat{\mu}/\alpha}_{\lambda/\beta,(m-1)}.
	\end{align*}
	It follows that
	\begin{align*}
		X &= -\sum_{j=1}^k \sum_{i=0}^{k-j}(-1)^{k-j+i}\sum_{\substack{\gamma\vdash k-j + (m-1)i\\ \alpha\vdash k+(m-1)i\\ \beta\vdash i}} c^{\alpha}_{\gamma,(1^j)}\cdot a^{\gamma/(k-j-i)}_{\beta',(m)}\cdot a^{\hat{\mu}/\alpha}_{\lambda/\beta,(m-1)}\\
		&= \sum_{i=0}^{k-1} \sum_{j=1}^{k-i} (-1)^{k+i} \cdot (-1)^{j+1} \sum_{\substack{\gamma\vdash k-j + (m-1)i\\ \alpha\vdash k+(m-1)i\\ \beta\vdash i}} c^{\alpha}_{\gamma,(1^j)}\cdot a^{\gamma/(k-j-i)}_{\beta',(m)}\cdot a^{\hat{\mu}/\alpha}_{\lambda/\beta,(m-1)}\\
		&= \sum_{i=0}^{k-1} (-1)^{k+i} \sum_{j=1}^{k-i} \sum_{\beta\vdash i} (-1)^{j+1}\cdot Y^\beta_{i,j},
	\end{align*}
	where we set
	\begin{align*}
		Y^\beta_{i,j}\coloneqq \sum_{\substack{\gamma\vdash k-j + (m-1)i\\ \alpha\vdash k+(m-1)i}} c^{\alpha}_{\gamma,(1^j)}\cdot a^{\gamma/(k-j-i)}_{\beta',(m)}\cdot a^{\hat{\mu}/\alpha}_{\lambda/\beta,(m-1)}.
	\end{align*}
	Next, we will simplify $Y^\beta_{i,j}$, for any fixed $\beta$, $i$ and $j$. To ease notation, in the rest of this proof we will abbreviate sums over all partitions of a given size. That is, we shorten $\sum_{\omega\vdash t}$ to $\sum_\omega$ (the size $t$ will always be clear from context). We use \eqref{eqn:skew-plethysm} to obtain
	\begin{align*}
		Y^\beta_{i,j}=\sum_{\gamma} \sum_\alpha c^\alpha_{\gamma,(1^j)}\cdot a^{\hat{\mu}/\alpha}_{\lambda/\beta,(m-1)} \cdot \sum_\varepsilon c^\gamma_{\varepsilon,(k-j-i)}\cdot a^\varepsilon_{\beta',(m)}.
	\end{align*}
	By \cite[(2.1)]{SLthesis}, we have that $\sum_\gamma c^\alpha_{\gamma,(1^j)}\cdot c^{\gamma}_{\varepsilon,(k-j-i)} = c^{\alpha}_{\varepsilon,(k-j-i),(1^j)} = \langle \chi^{\alpha/\varepsilon}, \chi^{k-j-i}\boxtimes \chi^{(1^j)}\rangle$. By Theorem~\ref{thm:LR}, $\chi^{k-j-i}\boxtimes \chi^{(1^j)} = \chi^{H(j)}+\chi^{H(j-1)}$ where $H(j)\coloneqq (k-i-j,1^j)$, except if $j=k-i$ then $\chi^\emptyset\boxtimes \chi^{(1^{k-i})}=\chi^{H(k-i-1)}$ only, i.e.~we treat the $\chi^{H(k-i)}$ term as the zero character. Hence
	\begin{align*}
		Y^\beta_{i,j} &= \sum_\alpha \sum_\varepsilon \langle \chi^{\alpha/\varepsilon}, \chi^{H(j)}+\chi^{H(j-1)}\rangle \cdot a^\varepsilon_{\beta',(m)}\cdot a^{\hat{\mu}/\alpha}_{\lambda/\beta,(m-1)}\\
		&= \sum_\alpha \sum_\varepsilon (c^\alpha_{\varepsilon,H(j)} + c^\alpha_{\varepsilon,H(j-1)}) \cdot a^\varepsilon_{\beta',(m)}\cdot a^{\hat{\mu}/\alpha}_{\lambda/\beta,(m-1)}
	\end{align*}
	(where we omit the $c^\alpha_{\varepsilon,H(k-i)}$ term). Since $H(0)=(k-i)$, we obtain
	\[ \sum_{j=1}^{k-i}(-1)^{j+1} \cdot Y^\beta_{i,j} = \sum_\alpha \sum_\varepsilon c^\alpha_{\varepsilon,(k-i)}\cdot a^\varepsilon_{\beta',(m)}\cdot a^{\hat{\mu}/\alpha}_{\lambda/\beta,(m-1)} = \sum_\alpha a^{\alpha/(k-i)}_{\beta',(m)}\cdot a^{\hat{\mu}/\alpha}_{\lambda/\beta,(m-1)}.\]
	Thus, we finally obtain
	\begin{align*} X &= \sum_{i=0}^{k-1}(-1)^{k+i} \cdot \sum_{\beta\vdash i} \sum_{j=1}^{k-i}(-1)^{j+1}\cdot Y^\beta_{i,j} \\&= \sum_{i=0}^{k-1}(-1)^{k+i}\cdot \sum_{\beta\vdash i}\ \sum_{\alpha\vdash k+(m-1)i} a^{\alpha/(k-i)}_{\beta',(m)}\cdot a^{\hat{\mu}/\alpha}_{\lambda/\beta,(m-1)}, \end{align*}
	as desired.
\end{proof}

To prove Theorem~\ref{thm:14.6ii}, we first deal with the case of $m=1$. In this case, $a^{\nu/(1^k)}_{\lambda',(1)} = c^{\nu}_{\lambda',(1^k)}$, which takes value 1 if $\nu-\lambda'$ is a sequence of 0s and 1s containing exactly $k$ many 1s, and takes value 0 otherwise. On the other hand, $\sum_{\alpha,\beta\vdash k} \big(a^\alpha_{\beta',(1)}\cdot a^{\hat{\nu}/\alpha}_{\lambda/\beta,\emptyset}\big)=c^{\lambda'}_{\hat{\nu},(1^{n-k})}$, which takes value 1 if $\lambda'-\hat{\nu}$ is a sequence of 0s and 1s containing exactly $n-k$ many 1s, and takes value 0 otherwise. We see that these two quantities are equal since $\nu=\hat{\nu}+(1^n)$, and hence Theorem~\ref{thm:14.6ii} holds when $m=1$. 

Next, we introduce some lemmas in preparation for proving Theorem~\ref{thm:14.6ii} when $m\ge 2$.

\begin{lemma}\label{lem:16.2}
	Let $r\in\N$ and $n_1,\dotsc,n_r\in\N_0$. Let $\gamma$ be a partition and for each $i\in\{1,2,\dotsc,r\}$, let $\psi_i$ be a virtual character of $S_{n_i}$. Then
	\[ (\psi_1\boxtimes\cdots\boxtimes\psi_r)/\chi^\gamma = \sum_{\gamma_1,\dotsc,\gamma_r} c^\gamma_{\gamma_1,\dotsc,\gamma_r}\cdot (\psi_1/\chi^{\gamma_1})\boxtimes\cdots\boxtimes(\psi_r/\chi^{\gamma_r}), \]
	summed over all sequences of partitions $\gamma_1,\dotsc,\gamma_r$ such that $|\gamma_1|+\cdots+|\gamma_r|=|\gamma|$.
\end{lemma}

\begin{proof}
	The case $r=1$ is trivial. For ease of notation we prove the statement for $r=2$; the case of general $r$ follows by an analogous argument. Define $n\coloneqq n_1+n_2$, $k\coloneqq |\gamma|$ and let $\delta\vdash n-k$. Then
	\begin{align*}
		\langle (\psi_1\boxtimes\psi_2)/\chi^\gamma,\chi^\delta\rangle &= \langle (\psi_1\times\psi_2)\up_{S_{n_1}\times S_{n_2}}^{S_n}, (\chi^\delta\times\chi^\gamma)\up_{S_{n-k}\times S_k}^{S_n}\rangle\\
		&= \langle (\psi_1\times\psi_2)\up_{S_{n_1}\times S_{n_2}}^{S_n}\down_{S_{n-k}\times S_k}, \chi^\delta\times\chi^\gamma\rangle,
	\end{align*}
	and applying Mackey's Theorem,
	\footnotesize
	\begin{align*}
		&= \sum_{\substack{0\le t_1,t_2\le k\\ t_1+t_2=k}} \langle (\psi_1\times\psi_2)\down^{S_{n_1}\times S_{n_2}}_{S_{n_1-t_1}\times S_{t_1}\times S_{n_2-t_2}\times S_{t_2}}\up^{S_{n-k}\times S_{k}}, \chi^\delta\times\chi^\gamma\rangle\\
		&= \sum_{t_1+t_2=k} \sum_{\substack{i=1,2\\\gamma_i\vdash t_i\\\delta_i\vdash n_i-t_i}} c^\gamma_{\gamma_1,\gamma_2}\cdot c^\delta_{\delta_1,\delta_2}\cdot \langle \psi_1\down_{S_{n_1-t_1}\times S_{t_1}} \times \psi_2\down_{S_{n_2-t_2}\times S_{t_2}}, \chi^{\delta_1}\times\chi^{\gamma_1}\times\chi^{\delta_2}\times\chi^{\gamma_2}\rangle\\
		&= \sum_{t_1+t_2=k} \sum_{\substack{i\\\gamma_i\vdash t_i\\\delta_i\vdash n_i-t_i}} c^\gamma_{\gamma_1,\gamma_2}\cdot c^\delta_{\delta_1,\delta_2}\cdot \langle \psi_1/\chi^{\gamma_1},\chi^{\delta_1}\rangle\cdot\langle \psi_2/\chi^{\gamma_2},\chi^{\delta_2}\rangle\\
		&= \sum_{t_1+t_2=k} \sum_{\substack{i\\\gamma_i\vdash t_i}} c^\gamma_{\gamma_1,\gamma_2}\cdot \langle (\psi_1/\chi^{\gamma_1})\boxtimes(\psi_2/\chi^{\gamma_2}), \chi^\delta\rangle,
	\end{align*}
	\normalsize
	as claimed.
\end{proof}

\begin{corollary}\label{cor:16.2}
	Let $m,n\in\N$ and $\lambda=(\lambda_1,\dotsc,\lambda_r)\vdash n$. Let $k\in\{0,1,\dotsc,n-1\}$ and $\beta\vdash k$. Then
	\begin{enumerate}[label=\textup{(\roman*)}]
		\item \label{591} $\big(\operatorname*{\boxtimes}_{i=1}^r \rho^{(\lambda_i)}_m \big)/\chi^{(1^{n-k})} = \sum_{\underline{t}}  \operatorname*{\boxtimes}_{i=1}^r \big( \rho^{(\lambda_i)}_m/\chi^{(1^{\lambda_i-t_i})} \big)$, and
		
		\item \label{592} $\big( \operatorname*{\boxtimes}_{i=1}^r \chi^{(1^{\lambda_i})} \big)/\chi^{\beta'} = \sum_{\underline{t}} c^\beta_{(t_1),\dotsc,(t_r)}\cdot \operatorname*{\boxtimes}_{i=1}^r \chi^{(1^{\lambda_i-t_i})}$,
	\end{enumerate}
	summed over all compositions $\underline{t}=(t_1,\dotsc,t_r)$ of $k$ into $r$ parts. That is, $t_i\in\N_0$ for all $i$ and $t_1+\cdots+t_r=k$ \textup{(}and we may further assume $t_i\le\lambda_i$ for all $i$\textup{)}.
\end{corollary}

\begin{proof}
	\ref{591} Applying Lemma~\ref{lem:16.2} with $\gamma=(1^{n-k})$, observe that $c^\gamma_{\gamma_1,\dotsc,\gamma_r}\in\{0,1\}$ and is non-zero only if each $\gamma_i=(1^{s_i})$ for some $s_i\in\N_0$. By Lemma~\ref{lem:tall-pleth}, $\rho^{(\lambda_i)}_m$ only has irreducible constituents $\chi^\mu$ where $l(\mu)\le \lambda_i$, so we may further assume that $s_i\le\lambda_i$. Writing $t_i=\lambda_i-s_i$ gives the result.
	
	\noindent \ref{592} Applying Lemma~\ref{lem:16.2} with $\psi_i=\chi^{(1^{\lambda_i})}$, observe that $\psi_i/\chi^{\gamma_i}\ne 0$ only if $\gamma_i=(1^{t_i})$ for some $0\le t_i\le \lambda_i$. Moreover, $c^{\beta'}_{(1^{t_1}),\dotsc,(1^{t_r})}=c^\beta_{(t_1),\dotsc,(t_r)}\in\{0,1\}$.
\end{proof}

\begin{lemma}\label{lem:16.3}
	Let $n\in\N$ and $k\in\{0,1,\dotsc,n\}$. Let $\nu$ be a partition with $l(\nu)=n$ and set $\hat{\nu}\coloneqq \nu-(1^n)$. Let $\delta\vdash|\nu|-k$ with $l(\delta)\le n$. Then
	\[ \langle\chi^\delta,\chi^{\nu/(1^k)} \rangle = \langle \chi^{\delta/(1^{n-k})},\chi^{\hat{\nu}}\rangle. \]
\end{lemma}

\begin{proof}
	The determinantal form of skew characters of symmetric groups (see e.g.~\cite[2.3.13]{JK}) gives $\chi^{\alpha/\beta} = \det\big(\chi^{(\alpha_i-i-\beta_j+j)}\big)$ whenever $\alpha$ and $\beta$ are partitions, where the multiplication of characters in expanding the determinant is given by the operation $\boxtimes$. Applying this to $\alpha=\nu'$ and $\beta=\emptyset$ and expanding the determinant with respect to the first row gives
	\[ \chi^{\nu'} = \sum_{j\ge 1} (-1)^{j-1}\cdot \chi^{(n+j-1)}\boxtimes \chi^{\hat{\nu}'/(1^{j-1})} = \sum_{j\ge 0} (-1)^j\cdot\chi^{(n+j)}\boxtimes \chi^{\hat{\nu}'/(1^j)}. \]
	Multiplying both sides by the sign representation then gives
	\[ \chi^\nu = \sum_{j\ge 0} (-1)^j\cdot \chi^{(1^{n+j})}\boxtimes \chi^{\hat{\nu}/(j)}. \]
	Then
	\begin{align*}
		\langle\chi^\delta,\chi^{\nu/(1^k)}\rangle &= \langle \chi^\delta\boxtimes\chi^{(1^k)},\chi^\nu\rangle \\
		&= \sum_{j\ge 0}(-1)^j\cdot\langle(\chi^\delta\boxtimes\chi^{(1^k)})/\chi^{(1^{n+j})}, \chi^{\hat{\nu}/(j)} \rangle\\
		&\overset{\normalfont\tiny \text{(Lemma~\ref{lem:16.2})}}{=} \sum_{j=0}^k (-1)^j\cdot \left\langle \sum_{s=0}^{k-j}(\chi^\delta/\chi^{(1^{n-s})})\boxtimes(\chi^{(1^k)}/\chi^{(1^{s+j})}), \chi^{\hat{\nu}/(j)}\right\rangle\\
		&=\sum_{j=0}^k \sum_{s=0}^{k-j}(-1)^j \langle (\chi^\delta/\chi^{(1^{n-s})})\boxtimes \chi^{(1^{k-j-s})}, \chi^{\hat{\nu}/(j)}\rangle\\
		&= \sum_{s=0}^k \left\langle (\chi^\delta/\chi^{(1^{n-s})})\boxtimes \left(\sum_{j=0}^{k-s}(-1)^j\cdot\chi^{(1^{k-s-j})}\boxtimes\chi^{(j)}\right), \chi^{\hat{\nu}}\right\rangle\\
		&=\langle \chi^\delta/\chi^{(1^{n-k})}, \chi^{\hat{\nu}}\rangle,
	\end{align*}
	where the final equality follows since $\chi^{(1^{k-s-j})}\boxtimes\chi^{(j)}=\chi^{(j+1,1^{k-s-j-1})}+\chi^{(j,1^{k-s-j})}$, and so $\sum_{j=0}^{k-s}(-1)^j\cdot\chi^{(1^{k-s-j})}\boxtimes\chi^{(j)}$ equals zero if $k\ne s$, and equals $\chi^\emptyset$ if $k=s$.
\end{proof}

\begin{lemma}\label{lem:16.6}
	Let $m$, $u$ and $t$ be integers with $m\ge 2$ and $u\ge t\ge 0$. Then
	\[ \rho^{(u)}_m/\chi^{(1^{u-t})} = \rho^{(t)}_m\boxtimes \rho^{(1^{u-t})}_{m-1}. \]
\end{lemma}

\begin{proof}
	Let $\delta\vdash mu-(u-t)$ be arbitrary. We show that $\langle \rho^{(u)}_m/\chi^{(1^{u-t})}, \chi^\delta\rangle = \langle\rho^{(t)}_m\boxtimes \rho^{(1^{u-t})}_{m-1}, \chi^\delta\rangle$. 
	Letting $H\coloneqq S_m\wr S_u$ and $K\coloneqq S_{|\delta|}\times S_{u-t}$, and substituting in the definition $\rho^{(u)}_m$ from Notation~\ref{not:rho}, we have by Mackey's theorem that
	\begin{align}\label{eqn:16.6}
		\langle \rho^{(u)}_m/\chi^{(1^{u-t})}, \chi^\delta\rangle &= \langle \rho^{(u)}_m, \chi^\delta\boxtimes\chi^{(1^{u-t})}\rangle = \langle \triv_H\up^{S_{mu}}, (\chi^\delta\times\chi^{(1^{u-t})})\up_K^{S_{mu}}\rangle \nonumber\\
		&= \sum_{\sigma\in K\setminus S_{mu}/H}\langle\triv_{H^\sigma}\down_{K\cap H^\sigma}\up^K,\chi^\delta\times\chi^{(1^{u-t})}\rangle\nonumber\\
		&=\sum_{\sigma\in K\setminus S_{mu}/H}\langle\triv_{K\cap H^\sigma},(\chi^\delta\times\chi^{(1^{u-t})})\down^K_{K\cap H^\sigma}\rangle,
	\end{align}
	where the final equality follows from Frobenius reciprocity.
	Here $\sigma$ runs over a set of representatives of double $(K,H)$-cosets in $S_{mu}$. Since $K\cap H^\sigma$ are the point stabilisers of the action of $K$ on the set of partitions of $\{1,2,\dotsc,mu\}$ into $u$ subsets of size $m$, the representatives $\sigma$ are parametrised by partitions of $u-t$ into exactly $u$ parts, including parts of size zero. 
	Fix one such partition $\sigma$ of $u-t$ and suppose that $\gamma_i$ is the number of parts of size $i$, for each $i\in\N_0$. Then $K\cap H^\sigma \cong \prod_{i\in\N_0} (S_{m-i}\times S_i)\wr S_{\gamma_i}$, and
	\begin{align*}
		\langle (\chi^\delta\times\chi^{(1^{u-t})})\down^K_{K\cap H^\sigma}, \triv\rangle &\le \langle (\chi^\delta\times\chi^{(1^{u-t})})\down_{\prod_i (S_{m-i}\times S_i)^{\times \gamma_i}}, \triv\rangle\\
		&= \langle \chi^\delta\down_{\prod_i S_{m-i}^{\times\gamma_i}},\triv\rangle\cdot \langle \chi^{(1^{u-t})}\down_{\prod_i S_i^{\times\gamma_i}},\triv\rangle.
	\end{align*}
	However, $\chi^{(1^{u-t})}$ is the sign representation, so $\langle \chi^{(1^{u-t})}\down_{\prod_i S_i^{\times\gamma_i}},\triv\rangle\ne 0$ if and only if $\gamma_i=0$ for all $i\ge 2$. Hence there is at most one $\sigma$ giving a non-zero contribution to the sum in \eqref{eqn:16.6}, namely $\sigma=(1^{u-t},0^t)$, and in this case $K\cap H^\sigma\cong (S_m\wr S_t)\times \big((S_{m-1}\times S_1)\wr S_{u-t}\big)$. Substituting into \eqref{eqn:16.6},
	\begin{align*}
		\langle \rho^{(u)}_m/\chi^{(1^{u-t})}, \chi^\delta\rangle &= \langle (\chi^\delta\times\chi^{(1^{u-t})})\down^K_{S_m\wr S_t\times (S_{m-1}\times S_1)\wr S_{u-t}},\ \triv\rangle,
		\end{align*}
which by Frobenius reciprocity equals
		\[\langle (\chi^\delta\times\chi^{(1^{u-t})})\down^{S_{|\delta|}\times S_{u-t}}_{S_m\wr S_t\times S_{m-1}\wr S_{u-t}\times S_1\wr S_{u-t}},\ \triv\up_{S_m\wr S_t\times (S_{m-1}\times S_1)\wr S_{u-t}}^{{S_m\wr S_t\times S_{m-1}\wr S_{u-t}\times S_1\wr S_{u-t}}}\rangle.
	\]
		 Noting that $|\delta|=mu-(u-t)=mt+(m-1)(u-t)$ and $\cX(\triv_{S_1};\chi^\omega)=\chi^\omega$, and using Lemma~\ref{lem:X-induced} in the second equality below, we have
	\begin{align*}
		\langle \rho^{(u)}_m/\chi^{(1^{u-t})}, \chi^\delta\rangle = \left\langle \chi^\delta\down^{S_{|\delta|}}_{S_m\wr S_t\times S_{m-1}\wr S_{u-t}} \times \chi^{(1^{u-t})},\ \triv_{S_m\wr S_t} \times \triv\up_{(S_{m-1}\times S_1)\wr S_{u-t}}^{S_{m-1}\wr S_{u-t}\times S_1\wr S_{u-t}}\right\rangle,
		\end{align*}
which in turn equals
\[
 \left\langle \chi^\delta\down^{S_{|\delta|}}_{S_m\wr S_t\times S_{m-1}\wr S_{u-t}} \times \chi^{(1^{u-t})},\ \triv_{S_m\wr S_t} \times \sum_{\omega\vdash u-t}\cX(\triv_{S_{m-1}}; \chi^\omega)\cdot\cX(\triv_{S_1};\chi^\omega) \right\rangle.
 \]
 This last expression simplifies to
\[
 \sum_{\omega\vdash u-t} \langle \chi^\delta\down_{S_m\wr S_t\times S_{m-1}\wr S_{u-t}},\ \triv_{S_m\wr S_t}\times \cX(\triv_{S_{m-1}};\chi^\omega)\rangle\cdot\langle \chi^{(1^{u-t})},\chi^\omega\rangle.
 \]
Now $\langle \chi^{(1^{u-t})},\chi^\omega\rangle=1$ precisely when $\omega=(1^{u-t})$ and is 0 otherwise, so
	\begin{align*} \langle \rho^{(u)}_m/\chi^{(1^{u-t})}, \chi^\delta\rangle &= \langle \chi^\delta\down_{S_m\wr S_t\times S_{m-1}\wr S_{u-t}},\ \triv_{S_m\wr S_t}\times \cX(\triv_{S_{m-1}};\chi^{(1^{u-t})})\rangle\\ &= \langle \chi^\delta,\ \rho^{(t)}_m\boxtimes \rho^{(1^{u-t})}_m\rangle \end{align*}
	by Frobenius reciprocity, recalling Notation~\ref{not:rho}\ref{562} and \ref{563}.
	Since $\delta$ was arbitrary, then $\rho^{(u)}_m/\chi^{(1^{u-t})} = \rho^{(t)}_m\boxtimes \rho^{(1^{u-t})}_{m-1}$ as desired.
\end{proof}

\begin{remark}
	When $m=2$, we can see from Proposition~\ref{prop:thrall} that $\rho^{(u)}_2/\chi^{(1^{u-t})} = \sum_\delta \chi^\delta = \rho^{(t)}_2\boxtimes\chi^{(1^{u-t})}$ where the sum is over all $\delta\vdash u+t$ with exactly $u-t$ many odd parts.%\hfill$\lozenge$
\end{remark}

Next, we generalise Lemma~\ref{lem:16.6} from the trivial partition $(u)$ to arbitrary partitions, giving Theorem B, after which it will be straightforward to deduce Theorem~\ref{thm:14.6ii}.

\begin{theorem}[Theorem B]\label{thm:16.7}
	Let $m,n\in\N$ with $m\ge 2$. Let $\lambda\vdash n$ and $k\in\{0,1,\dotsc,n-1\}$. Then
	\[ \rho^\lambda_m/\chi^{(1^{n-k})} = \sum_{\beta\vdash k} \rho^\beta_m\boxtimes \rho^{\lambda'/\beta'}_{m-1}. \]
\end{theorem}

\begin{proof}
	Following the notation for $\underline{t}$ in Corollary~\ref{cor:16.2} and letting $r=l(\lambda)$, observe that
	\begin{align*}
		\sum_{\gamma\vdash n} K_{\gamma,\lambda} &\cdot \big( \rho^\gamma_m/\chi^{(1^{n-k})} \big) = \Big( \sum_{\gamma\vdash n} K_{\gamma,\lambda}\cdot \rho^\gamma_m \Big)/\chi^{(1^{n-k})} \overset{\normalfont\tiny\text{\eqref{eqn:15.2}}}{=} \Big( \operatorname*{\boxtimes}_{i=1}^r \rho^{(\lambda_i)}_m\Big)/\chi^{(1^{n-k})}\\
		&\overset{\normalfont\tiny\text{(Corollary~\ref{cor:16.2}\ref{591})}}{=} \sum_{\underline{t}} \operatorname*{\boxtimes}_{i=1}^r \big( \rho^{(\lambda_i)}_m/\chi^{(1^{\lambda_i-t_i})} \big)\\
		&\overset{\normalfont\tiny\text{Lemma~\ref{lem:16.6}}}{=} \sum_{\underline{t}} \operatorname*{\boxtimes}_{i=1}^r \big( \rho^{(t_i)}_m \boxtimes \rho^{(1^{\lambda_i-t_i})}_{m-1} \big)\\
		&\overset{\normalfont\tiny\text{\eqref{eqn:15.2}}}{=} \sum_{\underline{t}} \sum_{\beta\vdash k} c^\beta_{(t_1),\dotsc,(t_r)}\cdot\rho^\beta_m \boxtimes \big(\operatorname*{\boxtimes}_{i=1}^r \rho^{(1^{\lambda_i-t_i})}_{m-1}\big)\\
		&\overset{\normalfont\tiny \text{Lemma~\ref{lem:cX}}}{=} \sum_{\beta\vdash k} \rho^\beta_m \boxtimes \cX\Big(\triv_{S_{m-1}}; \sum_{\underline{t}} c^\beta_{(t_1),\dotsc,(t_r)}\cdot \operatorname*{\boxtimes}_{i=1}^r \chi^{(1^{\lambda_i-t_i})} \Big)\up_{S_{m-1}\wr S_{n-k}}^{S_{(m-1)(n-k)}}\\
		&\overset{\normalfont\tiny\text{Corollary~\ref{cor:16.2}\ref{592}}}{=} \sum_{\beta\vdash k} \rho^\beta_m \boxtimes \cX \Big( \triv_{S_{m-1}}; \big( \operatorname*{\boxtimes}_{i=1}^r \chi^{(1^{\lambda_i})}\big) /\chi^{\beta'} \Big) \up_{S_{m-1}\wr S_{n-k}}^{S_{(m-1)(n-k)}}\\
		&\overset{\normalfont\tiny\text{\eqref{eqn:sign}}}{=} \sum_{\beta\vdash k} \rho^\beta_m \boxtimes \cX \Big( \triv_{S_{m-1}}; \big( \zeta^\lambda\cdot\sgn_{S_n} \big) /\chi^{\beta'} \Big) \up_{S_{m-1}\wr S_{n-k}}^{S_{(m-1)(n-k)}}\\
		&\overset{\normalfont\tiny\text{\eqref{eqn:sign} and \eqref{eqn:kostka}}}{=} \sum_{\beta\vdash k} \rho^\beta_m \boxtimes \cX \Big( \triv_{S_{m-1}}; \big( \sum_{\gamma\vdash n} K_{\gamma,\lambda}\cdot \chi^{\gamma'} \big) /\chi^{\beta'} \Big) \up_{S_{m-1}\wr S_{n-k}}^{S_{(m-1)(n-k)}} \\
		&= \sum_{\gamma\vdash n} K_{\gamma,\lambda} \cdot \Big( \sum_{\beta\vdash k} \rho^\beta_m \boxtimes \rho^{\gamma'/\beta'}_{m-1} \Big)\quad\text{since $\chi^{\gamma'}/\chi^{\beta'}=\chi^{\gamma'/\beta'}$ by Notation~\ref{not:rho}\ref{561}}.
	\end{align*}
	Since the matrix $(K_{\gamma,\lambda})_{\gamma,\lambda\vdash n}$ is invertible (in fact unitriangular if the partitions are ordered lexicographically, see e.g.~\cite[Chapter 2]{JK}), we deduce that $\rho^\gamma_m/\chi^{(1^{n-k})} = \sum_{\beta\vdash k} \rho^\beta_m \boxtimes \rho^{\gamma'/\beta'}_{m-1}$ for each $\gamma\vdash n$.
\end{proof}

\begin{lemma}\label{lem:cX}
	Let $m,r,a_1,\dotsc,a_r\in\N$, and let $n=\sum_{i=1}^r a_i$. For each $i\in\{1,\dotsc,r\}$, let $\nu_i\vdash a_i$. Then $\boxtimes_{i=1}^r \rho_m^{\nu_i} = \cX(\triv_{S_m};\boxtimes_{i=1}^r \chi^{\nu_i} )\up_{S_m\wr S_n}^{S_{mn}}$.
\end{lemma}

\begin{proof}
	The case $r=1$ follows from Notation~\ref{not:rho}\ref{562}. For each of notation we prove the statement for $r=2$; the case of general $r$ follows by an analogous argument. In fact, we can prove more generally that if $a,b\in\N$ and $\phi_1\in\Ch(S_a), \phi_2\in\Ch(S_b)$, then $\cX_{left}=\cX_{right}$ where
	\[ \cX_{left}\coloneqq  \left[\cX(\triv_{S_m}; \phi_1)\up_{S_m\wr S_a}^{S_{ma}} \times \cX(\triv_{S_m}; \phi_2)\up_{S_m\wr S_b}^{S_{mb}}\right]\up_{S_{ma}\times S_{mb}}^{S_{m(a+b)}} \]
	and
	\[ \cX_{right}\coloneqq  \cX\left(\triv_{S_m}; (\phi_1\times \phi_2)\up_{S_a\times S_b}^{S_{a+b}}\right)\up_{S_m\wr S_{a+b}}^{S_{m(a+b)}}, \]
	from which we recover the case of $r=2$ by setting $\phi_i=\chi^{\nu_i}$. 
	
	To prove that $\cX_{left}=\cX_{right}$, we first observe that $S_m\wr(S_a\times S_b) = S_m\wr S_a \times S_m\wr S_b$ (viewing $S_a\times S_b$ as a subgroup of $S_{a+b}$). Calling this group $U$, it is a subgroup of both $T_d\coloneqq S_{ma}\times S_{mb}$ and $T_w\coloneqq S_m\wr S_{a+b}$, and both $T_d$ and $T_w$ are subgroups of $S\coloneqq S_{m(a+b)}$. Now, by Lemma~\ref{lem:infl-ind} and the definition of $\cX(-;-)$,
	\begin{multline*}
		\cX(\triv_{S_m}; (\phi_1 \times \phi_2)\up_{S_a\times S_b}^{S_{a+b}}) = \cX(\triv_{S_m}; \phi_1\times\phi_2)\up_U^{T_w} = \left(\Infl_{S_a\times S_b}^{T_w}(\phi_1\times\phi_2)\right)\up_U^{T_w}\\
		= \left( \Infl_{S_a}^{S_m\wr S_a}(\phi_1)\times \Infl_{S_b}^{S_m\wr S_b}(\phi_2) \right)\up_U^{T_w} = \Big( \cX(\triv_{S_m};\phi_1)\times \cX(\triv_{S_m};\phi_2) \Big)\up_U^{T_w}.
	\end{multline*}
	Therefore
	\begin{multline*}
	 	\cX_{right} = \Big( \cX(\triv_{S_m};\phi_1)\times \cX(\triv_{S_m};\phi_2) \Big)\up_U^{T_w} \up_{T_w}^S \\
	 	= \Big( \cX(\triv_{S_m};\phi_1)\times \cX(\triv_{S_m};\phi_2) \Big)\up_U^{T_d} \up_{T_d}^S = \cX_{left},
	 \end{multline*}
	where the second equality follows from the transitivity of induction.
\end{proof}

\begin{proof}[Proof of Theorem~\ref{thm:14.6ii} when $m\ge 2$]
	Take $\langle -, \chi^{\hat{\nu}}\rangle$ in Theorem~\ref{thm:16.7} to obtain
	\[ \langle \rho^\lambda_m/\chi^{(1^{n-k})}, \chi^{\hat{\nu}} \rangle = \left\langle \sum_{\beta\vdash k} \rho^\beta_m\boxtimes \rho^{\lambda'/\beta'}_{m-1}, \chi^{\hat{\nu}} \right\rangle. \]
	By Lemma~\ref{lem:16.3}, $\langle \rho^\lambda_m/\chi^{(1^{n-k})}, \chi^{\hat{\nu}} \rangle = \langle \rho^\lambda_m, \chi^{\nu/(1^k)}\rangle = a^{\nu/(1^k)}_{\lambda,(m)}$. On the other hand,
	\[ \left\langle \sum_{\beta\vdash k} \rho^\beta_m\boxtimes \rho^{\lambda'/\beta'}_{m-1}, \chi^{\hat{\nu}} \right\rangle = \sum_{\beta\vdash k} \sum_{\alpha\vdash mk} \langle \rho^\beta_m,\chi^\alpha\rangle \cdot \langle \rho^{\lambda'/\beta'}_{m-1}, \chi^{\hat{\nu}/\alpha}\rangle = \sum_{\substack{\alpha\vdash mk\\\beta\vdash k}} a^\alpha_{\beta,(m)}\cdot a^{\hat{\nu}/\alpha}_{\lambda'/\beta',(m-1)}, \]
	which concludes the proof.
\end{proof}
% END NATIVE AUTHOR BLOCK

\begin{thebibliography}{99}
\bibitem{BBP}Bessenrodt, C.; Bowman, C.; Paget, R. The classification of multiplicity-free plethysms of Schur functions, Trans. Amer. Math. Soc., Volume 375 (2022) no. 7, pp. 5151-5194
\bibitem{BOR}Briand, E.; Orellana, R.; Rosas, M. Rectangular symmetries for coefficients of symmetric functions, Electron. J. Combin., Volume 22 (2015) no. 3, p. Paper 3.15, 18 pages
\bibitem{Brion}Brion, M. Stable properties of plethysm: on two conjectures of Foulkes, Manuscripta Math., Volume 80 (1993) no. 4, pp. 347-371
\bibitem{BCV}Bruns, W.; Conca, A.; Varbaro, M. Relations between the minors of a generic matrix, Adv. Math., Volume 244 (2013), pp. 171-206
\bibitem{BCI}Bürgisser, P.; Christandl, M.; Ikenmeyer, C. Even partitions in plethysms, J. Algebra, Volume 328 (2011), pp. 322-329
\bibitem{dB-thesis}de Boeck, M. On the structure of Foulkes modules for the symmetric group, Ph. D. Thesis, University of Kent (2015)
\bibitem{dBPW}de Boeck, M.; Paget, R.; Wildon, M. Plethysms of symmetric functions and highest weight representations, Trans. Amer. Math. Soc., Volume 374 (2021) no. 11, pp. 8013-8043
\bibitem{EL}Erdös, P.; Lehner, J. The distribution of the number of summands in the partitions of a positive integer, Duke Math. J., Volume 8 (1941), pp. 335-345
\bibitem{EPW}Evseev, A.; Paget, R.; Wildon, M. Character deflations and a generalization of the Murnaghan-Nakayama rule, J. Group Theory, Volume 17 (2014) no. 6, pp. 1035-1070
\bibitem{Foulkes}Foulkes, H. O. Concomitants of the quintic and sextic up to degree four in the coefficients of the ground form, J. London Math. Soc., Volume 25 (1950), pp. 205-209
\bibitem{G}Giannelli, E. Characters of odd degree of symmetric groups, J. Lond. Math. Soc. (2), Volume 96 (2017) no. 1, pp. 1-14
\bibitem{GKNT}Giannelli, E.; Kleshchev, A.; Navarro, G.; Tiep, P. H. Restriction of odd degree characters and natural correspondences, Int. Math. Res. Not. IMRN (2017) no. 20, pp. 6089-6118
\bibitem{GL1}Giannelli, E.; Law, S. On permutation characters and Sylow p-subgroups of S n , J. Algebra, Volume 506 (2018), pp. 409-428
\bibitem{GL2}Giannelli, E.; Law, S. Sylow branching coefficients for symmetric groups, J. Lond. Math. Soc. (2), Volume 103 (2021) no. 2, pp. 697-728
\bibitem{GLLV}Giannelli, E.; Law, S.; Long, J.; Vallejo, C. Sylow branching coefficients and a conjecture of Malle and Navarro, Bull. Lond. Math. Soc., Volume 54 (2022) no. 2, pp. 552-567
\bibitem{GN}Giannelli, E.; Navarro, G. Restricting irreducible characters to Sylow p-subgroups, Proc. Amer. Math. Soc., Volume 146 (2018) no. 5, pp. 1963-1976
\bibitem{I}Isaacs, I. M. Character theory of finite groups, Dover Publications, Inc., New York, 1994, xii+303 pages
\bibitem{INOT}Isaacs, I. M.; Navarro, G.; Olsson, J. B.; Tiep, P. H. Character restrictions and multiplicities in symmetric groups, J. Algebra, Volume 478 (2017), pp. 271-282
\bibitem{JK}James, G.; Kerber, A. The representation theory of the symmetric group, Encyclopedia of Mathematics and its Applications, 16, Addison-Wesley Publishing Co., Reading, Mass., 1981, xxviii+510 pages
\bibitem{J}James, G. D. The representation theory of the symmetric groups, Lecture Notes in Mathematics, 682, Springer, Berlin, 1978, v+156 pages
\bibitem{LR}Langley, T. M.; Remmel, J. B. The plethysm s \(\lambda\) [s \(\mu\) ] at hook and near-hook shapes, Electron. J. Combin., Volume 11 (2004) no. 1, p. Research Paper 11, 26 pages
\bibitem{SLthesis}Law, S. On Problems in the Representation Theory of Symmetric Groups, Ph. D. Thesis, University of Cambridge (2019)
\bibitem{Mac}Macdonald, I. G. Symmetric functions and Hall polynomials, Oxford Mathematical Monographs, The Clarendon Press, Oxford University Press, New York, 1995, x+475 pages (With contributions by A. Zelevinsky, Oxford Science Publications)
\bibitem{Man}Manivel, L. Gaussian maps and plethysm, Algebraic geometry (Catania, 1993/Barcelona, 1994) (Lecture Notes in Pure and Appl. Math.), Volume 200, Dekker, New York, 1998, pp. 91-117
\bibitem{MM}Manivel, L.; Michałek, M. Effective constructions in plethysms and Weintraub’s conjecture, Algebr. Represent. Theory, Volume 17 (2014) no. 2, pp. 433-443
\bibitem{N}Navarro, G. Character tables and Sylow subgroups revisited, Group theory and computation (Indian Stat. Inst. Ser.), Springer, Singapore, 2018, pp. 197-206
\bibitem{Olsson}Olsson, J. B. Combinatorics and representations of finite groups, Vorlesungen aus dem Fachbereich Mathematik der Universität GH Essen [Lecture Notes in Mathematics at the University of Essen], 20, Universität Essen, Fachbereich Mathematik, Essen, 1993, ii+94 pages
\bibitem{PW}Paget, R.; Wildon, M. Generalized Foulkes modules and maximal and minimal constituents of plethysms of Schur functions, Proc. Lond. Math. Soc. (3), Volume 118 (2019) no. 5, pp. 1153-1187
\bibitem{Stanley}Stanley, R. P. Enumerative combinatorics. Vol. 2. With a foreword by Gian-Carlo Rota and appendix 1 by Sergey Fomin, Cambridge Studies in Advanced Mathematics, 62, Cambridge University Press, Cambridge, 1999, xii+581 pages
\bibitem{Sta-OP}Stanley, R. P. Positivity problems and conjectures in algebraic combinatorics, Mathematics: frontiers and perspectives, Amer. Math. Soc., Providence, RI, 2000, pp. 295-319
\bibitem{T}Thrall, R. M. On symmetrized Kronecker powers and the structure of the free Lie ring, Amer. J. Math., Volume 64 (1942), pp. 371-388
\bibitem{W}Weintraub, S. H. Some observations on plethysms, J. Algebra, Volume 129 (1990) no. 1, pp. 103-114
\end{thebibliography}
\end{document}
